\section{Taylor--Wiles tests on a chain of curve points}
\label{sec:marked-patching}
\hypertarget{source-ordinary-marked-patching}{}

To move between components of ordinary Hecke support, we shall compare a chain of
residual curve points at one marked level. Choosing Taylor--Wiles primes separately at
each point would leave us without a common comparison.
We work with the normalized representation over the
characteristic-\(\ell\) power-series ring and test its
off-diagonal entries and local flags. The tests give the
patching hypotheses used in the component comparison of
Appendix~\ref{sec:linked-components}. Lemma~\ref{hyp:marked-square-zero-data} fixes the curve representation, its cyclotomic kernel and its coefficient field before any test elements or auxiliary places are chosen.

\begin{established}[Integral representatives at a normalized ordinary curve]\label{hyp:marked-square-zero-data}
Use a fitted point of Proposition~\ref{prop:fitted-points-three-tests}
at one of the three marked levels. Let \(G\) denote the CM subgroup,
\(H=\ker(\epsilon|_G)\) its full cyclotomic kernel, \(c\) the chosen
complex conjugation and \(\sigma h=chc^{-1}\). Let \(A=K[[t]]\)
and \(E=K((t))\) at the normalized curve, with \(K\) finite of
odd characteristic \(\ell\). In the polarized coordinates the
curve representation satisfies \(r(\sigma h)=\mu(h)r(h)^{-\mathsf T}\),
where \(\mu\) is the level multiplier. The span \(S=\operatorname{span}_E r(H)\) and the off-diagonal module
\(M\) are
\[
 (S,M)=\begin{cases}
 (M_n(E),0),&\text{type I},\\
 (M_a(E)\oplus M_b(E),\text{the off-diagonal blocks}),
      &\text{type II},\quad a+b=n .
 \end{cases}
\]

Use the conjugation action on \(M\) and \(-\mathsf T\) at \(c\).
The subspace of \(H^1\) whose classes have an upper triangular
representative in each ordered ordinary flag and are unramified on
the prescribed inertia away from \(\ell\) has a finite basis
represented by integral cocycles \(\kappa_1,\ldots,\kappa_g\) with
\[
 \kappa_j(c)=0,\qquad
 \kappa_j(\sigma h)=-\kappa_j(h)^{\mathsf T}.
\]
One can choose the representatives and their local changes of
frame over a single denominator.
\end{established}
\begin{proof}
The two span formulas follow from the irreducible, pairwise
inequivalent constituents on \(H\) in
Proposition~\ref{prop:fitted-points-three-tests}. In type II, if \(V_1,V_2\) denote these constituents, their matrix
coefficient modules give
\[
 E[r(H)]=\operatorname{End}_E(V_1)\oplus\operatorname{End}_E(V_2),
 \qquad M=\operatorname{Hom}_E(V_1,V_2)\oplus\operatorname{Hom}_E(V_2,V_1).
\]

The marked
auxiliary place in Proposition~\ref{prop:separated-marked-host}
supplies an element of \(G\) with nontrivial residual cyclotomic
character. For the finitely ramified real and CM groups and their
open subgroups \(U\), class field theory gives
\(\#\operatorname{Hom}_{\rm cont}(U,\mathbf F_\ell)<\infty\):
only the finite class group, the units at finitely many finite
places and finitely many uniformizers contribute modulo \(\ell\).
After passing to an open subgroup acting trivially on the residual
lattice, the action and integral cocycles factor through its
maximal pro-\(\ell\) quotient. 
For that quotient \(P=U^{(\ell)}\), finite generation follows from
\[
      \dim_{\mathbf F_\ell}P/\overline{P^\ell[P,P]}
       =\dim_{\mathbf F_\ell}\operatorname{Hom}_{\rm cont}(U,\mathbf F_\ell)<\infty .
\] Evaluating cocycles on
these generators and a finite set of coset representatives
bounds the dimension of the full cohomology, hence of the
subspace cut out by the local conditions.

For any cocycle \(z\), the identity \(c^2=1\) gives
\(c z(c)=-z(c)\). Subtract the coboundary of \(-z(c)/2\) to
obtain a representative vanishing at \(c\). The cocycle identity
at \(chc^{-1}\) then gives the transpose formula. Its continuous
image on the compact Galois group is bounded. Choose such
representatives for a basis, together with the finitely many
local changes permitted by the definition of the ordinary
subspace. A common power of \(t\) clears their denominators
without changing independence over \(E\). Boundaries between
normalized representatives have value zero at \(c\), so their
defining matrices satisfy \(Y=-Y^{\mathsf T}\).
\end{proof}

\subsection{The paired determinant test}
\label{subsec:det-test}

Begin at one curve point. The determinant must detect the cocycles together with the
underlying representation, so over \(E\) we form the finite algebra
\[
 B=S\oplus M^g,\qquad
 (x,p)(y,q)=(xy,xq+py).
\]
The second summand has square zero. The homomorphism
\[
 R(h)=\bigl(r(h),(\kappa_j(h)r(h))_j\bigr)
\]
takes values in \(B^\times\) by the cocycle rule. Entrywise transpose
on each component is an anti-involution \(\tau\) of \(B\), and
\(R(\sigma h)=\mu(h)\tau(R(h))^{-1}\). All these formulas are defined
on the evident free \(A\)-lattices.

\begin{lemma}[Spanning the marked square-zero algebra]\label{lem:marked-square-zero-span}
For the normalized cocycles and block decomposition of
Lemma~\ref{hyp:marked-square-zero-data}, one has
\[
   B=\operatorname{span}_E\{R(h):h\in H\}.
\]
\end{lemma}

\begin{proof}
\proofstep{Projection onto the two radical directions.}
The span of \(R(H)\) is a subalgebra
mapping onto \(S\). Its intersection \(W\) with \(M^g\) is an
\(S\)-bimodule, since the ideal has square zero. In the one-block
case there is no second summand. In the two-block case write
\(M=M_{12}\oplus M_{21}\). These are distinct simple split
\(S\)-bimodules. The central idempotents separate directions. The matrix units in each block then recover all of a submodule from its coefficient subspace:
\[
 W_{ij}=M_{ij}\otimes_E U_{ij}\subseteq M_{ij}\otimes_E E^g,\qquad U_{ij}\subseteq E^g .
\] If \(W\) is proper, take a nonzero row \(d=(d_j)\) annihilating
one of the \(U_{ij}\). For a lift \((x,(p_j))\) in the span of
\(R(H)\), define \(D(x)=\sum_j d_j p_j^{ij}\). Changing the lift
adds an element of \(M_{ij}\otimes U_{ij}\), which this row kills.
The square-zero multiplication rule gives
\[
                  D(xy)=xD(y)+D(x)y.
\]
Thus \(D\) is an \(E\)-linear derivation into the \(ij\)-corner.
\proofstep{The coboundary forced by a missing direction.}
For \(e_i\) the central idempotent of
the target block, put \(Y_{ij}=D(e_i)\). On the \(ij\)-corner the left action of \(e_i\) is one
and its right action is zero. Differentiating \(e_i x=xe_i\)
therefore gives
\[
                  Y_{ij}x + D(x) = xY_{ij}.
\] It follows
that the \(ij\)-part of \(\sum_j d_j\kappa_j(h)\), for \(h\in H\),
is \(r(h)Y_{ij}r(h)^{-1}-Y_{ij}\). Polarization at \(\sigma h\)
gives the other part with matrix \(-Y_{ij}^{\mathsf T}\). The
difference from this skew coboundary vanishes on \(H\). Its
values on \(G\) are then \(H\)-invariant by normality, and an
off-diagonal matrix commuting with \(r(H)\) is zero by the
two block idempotents. At \(c\) the difference is zero by
normalization and skew symmetry. This contradicts independence
of the \(\kappa_j\) and proves the spanning assertion.

\end{proof}

\paragraph{The twisted dual and the trace test.}
To read the determinant test as a condition on a linear functional, identify \(B^\vee=\operatorname{Hom}_E(B,E)\) with pairs
\(z=(\phi,(\lambda_j)_j)\) by the perfect pairing
\[
 \langle z,(x,p)\rangle=\operatorname{tr}(\phi x)
     +\sum_j\operatorname{tr}(\lambda_jp_j),
 \qquad \phi\in S,\quad\lambda_j\in M .
\]
The dual bimodule acts by \((a z b)(d)=z(bda)\).
Thus the trace pairing exchanges the two off-diagonal
directions. In this pairing \(\phi\) has coefficient \(\eta\), and \(\lambda_j\) has coefficient
\(\varepsilon_j'\), paired with \(\varepsilon_j\) by
\(\varepsilon_i\varepsilon_j'=\delta_{ij}\eta\).
The action on this dual is conjugation multiplied by the scalar twist:
\[
 h\cdot z=\chi(h)R(h)zR(h)^{-1}.
\]

In evaluating this action, conjugate the argument \((x,p)\)
of the functional by \(R(h)^{-1}\). Factoring
\(R(h)=(1,\kappa(h))(r(h),0)\) shows that the radical
part of the argument acquires the commutator
\(x\kappa_j(h)-\kappa_j(h)x\) before conjugation by
\(r(h)^{-1}\). Moving it across the trace gives the
commutator in the first coordinate below. The second coordinates are
\(\chi(h)r(h)\lambda_jr(h)^{-1}\). Its first coordinate is
\[
 \chi(h)\left(r(h)\phi r(h)^{-1}
       +\sum_j[\kappa_j(h),r(h)\lambda_jr(h)^{-1}]\right).
\]
Here \([P,Q]=PQ-QP\). Conjugation at \(c\) acts as
\(s\tau^*\), where \(\tau^*\) is transpose under the pairing
and \(s=-\chi(c)\). In particular \(s^2=1\).
For a twisted cocycle with \(z(c)=0\), the test at \(h\in H\) is
\[
 \mathcal L_h(z)=\langle z(h),R(h)\rangle
  =\operatorname{tr}(\phi(h)r(h))
     +\sum_j\operatorname{tr}(\lambda_j(h)\kappa_j(h)r(h)).
 \tag{10}\label{eq:local-patch-10}
\]
This is the \(\eta\)-coefficient of the infinitesimal trace.

\begin{lemma}[Paired trace detection]\label{lem:paired-trace-detection}
With the curve and block data of Lemma~\ref{hyp:marked-square-zero-data},
let \(z\) be a cocycle with values in the \(\chi\)-twisted dual and
\(z(c)=0\). The detection implication over \(E\) is
\[
  (\forall h\in H,\ \mathcal L_h(z)=0)\quad\Longrightarrow\quad [z]=0
\]
in first cohomology with these coefficients. For an integral cocycle
one can clear the boundary matrix by a power of \(t\).
\end{lemma}

\begin{proof}
\proofstep{From the vanishing trace to a derivation.}
Suppose all the tests vanish. For \(h\in H\), set
\(D_h=z(h)R(h)\in B^\vee\). Then
\(D_{hh'}=D_hR(h')+R(h)D_{h'}\).
If \(\sum_i a_iR(h_i)=0\), evaluate \(\sum_i a_i D_{h_i h'}\)
at the identity of \(B\). The vanishing of the tests and the
relation imply
\[
       \left\langle\sum_i a_i D_{h_i},R(h')\right\rangle=0 .
\]
Lemma~\ref{lem:marked-square-zero-span} gives
\(\sum_i a_i D_{h_i}=0\).
There is consequently a derivation \(D:B\to B^\vee\)
with \(D(R(h))=D_h\).

\proofstep{The action of the cyclotomic quotient.}
For \(v\in G\), put \(u_v=R(v)\) and \(w_v=D(u_v)u_v^{-1}\).
Write \(\delta_Y(x)=xY-Yx\). Applying the derivation rule to
\(u_v x u_v^{-1}\) gives
\[
 u_v D(u_v^{-1}x u_v)u_v^{-1}=D(x)+\delta_{w_v}(x).
\]
The cocycle rule on \(v h v^{-1}\), with \(h\in H\), gives
\[
 D(x)=\chi(v)u_v D(u_v^{-1}x u_v)u_v^{-1}
                                  -\delta_{z(v)}(x)
 \quad (x=R(v h v^{-1})).
\]
By linearity this holds on \(B\). Choosing \(v\) with
\(\chi(v)\ne1\) shows that \(D\) is inner. Subtract its twisted
coboundary from \(z\). The resulting cocycle vanishes on \(H\),
and takes values in the invariant subspace of \(B^\vee\) under
\(R(H)\). Such a functional satisfies \(xY=Yx\) for every
\(x\in B\) by the span.

On this subspace the action of \(G\)
is therefore the scalar \(\chi\). Since \(G/H\) is abelian,
the cocycle \(z_0\) on it satisfies
\[
 (\chi(v)-1)z_0(v')=(\chi(v')-1)z_0(v).
\]
The same choice of \(v\) makes \(z_0\) a coboundary. In the final
boundary matrix replace \(Y\) by \((Y+s\tau^*(Y))/2\).
The polarized cocycle rule and the formula for \(R(\sigma h)\)
imply that the new matrix gives the same boundary on \(G\),
and it is fixed by the action at \(c\). Thus it also gives the
boundary on the full real group. This proves the test after
inverting \(t\), or, integrally, after multiplying the cocycle
by a power of \(t\) to clear the matrix denominators.

\end{proof}

\begin{construction}
\[
\begin{array}{rl}
1&\text{Fix the finite chain of marked curve points.}\\
2&\text{Choose one finite set of tests in the cyclotomic kernel.}\\
3&\text{Fix the coefficient depth }m.\\
4&\text{At each point test at the base and at } N\gg_m1.\\ \hline
 &\#\{\text{tested dual classes at the point}\}\le |K|^{C},
      \quad C\text{ independent of }m .
\end{array}
\]
\caption{Order of choices for the paired determinant tests. The finite test set works across the fixed chain, while the bound on size may depend on the point.}\label{fig:marked-test-quantifiers}
\end{construction}

\subsection{Common tests and the finite-depth dual bound}

At a single point we have a trace test. Its use along a chain requires two kinds of
uniformity. The tests must be chosen in common, and their error at each point must
stay bounded as the torsion depth grows.

\begin{proposition}[Uniform bounds for the tested dual classes]
\label{prop:marked-uniform-dual-tests}
For the fixed finite chain, a common finite choice of tests in the
cyclotomic kernel detects the stated paired dual classes. For each
point there is a constant \(C\) independent of \(m\) so that the tested dual cohomology group \(H_{\rm test}\) at depth \(m\) satisfies 
\[
 |H_{\rm test}|\le|K|^C\qquad\text{at the base and at every auxiliary level }N\gg_m1.
\]
 The finite set
of tests is chosen before the auxiliary levels.
\end{proposition}

\begin{proof}
\proofstep{A fixed generating set for the finite-depth cocycles.}
For the standard adjoint module, the odd module and the paired
dual, the action on an open normal subgroup has pro-\(\ell\)
image. For these finite-free \(A\)-modules one can take the
kernel of the reduction of the action. At all finite
truncations the kernel of reduction consists of successive
additive matrix groups. At depth \(m\) write \(V_m\) for the additive coefficient module and \(\rho_m\) for its action. The cocycle equation is precisely the multiplication rule 
\[
 (z(xy),\rho_m(xy))=(z(x)+\rho_m(x)z(y),\rho_m(x)\rho_m(y)).
\]
 The image on the chosen subgroup therefore lies in an extension of two finite \(\ell\)-groups. The maximal pro-\(\ell\) quotient of the
open subgroup has finitely generated Frattini quotient by the
stated Galois finiteness, and hence is topologically finitely
generated. Adding lifts of the finite quotient gives a fixed
finite generating set on which all the cocycles are
determined, integrally and modulo every \(t^m\).

\proofstep{Finite test functionals on the common dense set.}
On a finite chain of points, regular semisimplicity on \(H\) is
an open dense condition for each point by the nice-point
hypotheses. Its openness follows by reducing the nonzero
discriminant at sufficiently high depth. The intersection is
dense, and by continuity a cocycle vanishing under \eqref{eq:local-patch-10} there
passes every test on \(H\). Evaluation at a fixed finite generating set \(g_1,\ldots,g_d\) identifies a cocycle by its values. For the coefficient space \(V\) at this point the injection is
\[
 Z^1(G,V)\hookrightarrow V^d,\qquad z\longmapsto(z(g_1),\ldots,z(g_d)).
\] Finitely many of the tests span the same space of
linear forms as the tests on the entire intersection. Take
the union of these finite sets for all points, say
\(h_1,\ldots,h_q\).

Split the characteristic polynomials of these elements by
finite extensions of each curve field, using the integer
rings for the extensions. The support kernels are unchanged
because the ring maps are injective. The rings are complete
equicharacteristic DVRs with finite residue fields. They can
be enlarged again until the residue fields all have the same
cardinality, as follows. If their degrees over \(\mathbf F_\ell\) are \(d_1,\ldots,d_s\), take unramified extensions to the field of cardinality 
\[
 \ell^{\operatorname{lcm}(d_1,\ldots,d_s)}.
\]
 Each coefficient extension is finite
free over the preceding DVR.

For a finite free extension \(A\subset A'\) of these
DVRs, choose a basis \(b_i\) of \(A'\) over \(A\).
An \(A'\)-valued cochain is a unique sum \(z=\sum_i b_i z_i\)
with continuous \(A\)-valued cochains \(z_i\). For an \(A\)-linear
equation \(\mathcal T\) in the cocycle and the matrices imposing its local conditions,
\[
       \mathcal T(z)=\sum_i b_i \mathcal T(z_i).
\]
Hence the equations cut out the scalar extensions of their
kernels. The image of the boundary map is also the scalar
extension of its image, and so is the cohomology quotient.
After inverting the uniformizer this preserves the chosen odd
cohomology basis. In particular the new tests still give precisely
the coboundaries over the new fraction fields.

\proofstep{Frobenius approximation at every fixed depth.}
At depth \(N\) use the finite image of the product
of the current marking quotient, the CM sign, and
\(\epsilon_\ell\bmod\ell^N\). The image of \(h_j\) is
\[
            (h_j\bmod V_N,\ 1,\ 1),
\]
since \(h_j\) is in the positive full cyclotomic kernel.
Thus the simultaneous congruence is realized in that
finite quotient. Chebotarev for its inverse class gives
geometric Frobenius in the desired class. Excluding the
fixed level and previously chosen primes still leaves a
choice of distinct split places. The cyclotomic coordinate specifies the norm of every chosen place: 
\[
 Q_{j,N}\equiv1\pmod{\ell^N}.
\]
 The auxiliary
level retains the chosen decomposition groups and tame
inertia generators.

Passing to the open subgroup cutting
out the residual diagonal characters and the mod-\(\ell\)
cyclotomic character, the point representations have
unipotent residual image. A finite group consisting of
unipotent matrices is an \(\ell\)-group by Cauchy's theorem.
Its congruence kernels and the cocycle semidirect images
are likewise \(\ell\)-groups at finite depth. This is the
factorization through the pro-\(\ell\) quotient used in
comparing the level kernels. Fix \(m\). The coefficient actions and base cocycles at precision \(t^m\) are recorded by finitely many maps to finite groups. If \(N_i\) is the open normal kernel for the \(i\)-th such map, they all factor through
\[
 G/\bigcap_i N_i\ \hookrightarrow\ \prod_i G/N_i .
\] Frobenius agreement thus holds eventually
for all of them at once.

\proofstep{The Smith bounds for the tested complex.}
Here is the integral size comparison. Use evaluation on the
finite generating set to write the entire cocycle and test
system as linear equations on \(A^d\), where \(d\) here denotes
the number of scalar generator coordinates. On a free
profinite group the tuples define cocycles by the
semidirect product universal property. The relations of its
quotient, the condition at \(c\) and \eqref{eq:local-patch-10} give rows over
\(A\). The submodule of the dual of \(A^d\) spanned by the
rows is finitely generated, so the same finite matrix
\(F\) gives the system at every truncation.

The matrix
\(\Delta\) of the coboundaries from the \(c\)-fixed subspace
has image in its kernel. That subspace commutes with
reduction since the eigenprojections of the involution
use division by two. Over the fraction field, passing the tests is equivalent to being a boundary. Thus 
\[
 \operatorname{im}(\Delta\otimes_A E)=\ker(F\otimes_A E).
\]
 If their common dimension is \(j\), Smith
normal form gives
\[
        |\ker F_m|\le |K|^{mj+C_1},\qquad
        |\operatorname{im}\Delta_m|\ge |K|^{mj-C_2}
\]
with constants given by valuations of the nonzero
elementary divisors. This bounds the size of the tested
cohomology by \(|K|^{C_1+C_2}\) for all \(m\).

\proofstep{Descent of the tested auxiliary classes.}
A dual class at an auxiliary level is tested if it
restricts to a coboundary on each new decomposition
group. In that case its value at tame inertia is zero,
because the coefficient action is pulled back from the
base level. Its cocycle descends to the base: the
auxiliary kernel is generated, after passing to the
pro-\(\ell\) quotient, by the normal closure of those tame
generators. On the auxiliary kernel the action on the coefficients is trivial. Thus a cocycle that vanishes on a generator vanishes on its conjugates by the cocycle rule, and then on the closed subgroup they generate by continuity:
\[
 z(ghg^{-1})=g\,z(h)\qquad(h\text{ in the kernel}).
\] Normalize
the cocycle at \(c\) by subtracting the coboundary of
\(-z(c)/2\).

At each new Frobenius the normalized value
is still a local coboundary, whose product trace against
the Frobenius matrix vanishes by cyclicity. At sufficiently
large \(N\) the agreement modulo \(t^m\) transports this
vanishing to \eqref{eq:local-patch-10} on the chosen \(h_j\). Boundaries descend
with it, so the bound for the tested cohomology applies
at the auxiliary level. The dual bound is uniform in the sense \(\exists C\,\forall m\,\forall N\gg_m 1\). The constant is
allowed to depend on the fixed curve point and finite tests.

\end{proof}

\subsection{The first-order bound}

We next apply the common tests to first-order deformations. Passing from normalized
cocycles to their quotient by boundaries is essential for the dimension count in the
completed local ring.

\begin{proposition}[First-order cohomology at a marked point]
\label{prop:marked-first-order-bound}
With the standard adjoint and chosen local conditions at a marked
point, let \(Z^1_{m,\mathcal L}\) be the normalized cocycles at depth \(m\) satisfying these conditions and \(B_m^1\) their normalized coboundaries. Then
\[
                       \bigl|Z^1_{m,\mathcal L}/B_m^1\bigr|\le |K|^{m(nd+nq)+C},
\]
where \(d=[L^+:\mathbb Q]\) and \(q\) is the number of auxiliary
places in this comparison. The constant is independent of \(m\)
and the bound holds on the shared auxiliary system for \(N\gg_m1\).
\end{proposition}

\begin{proof}
\proofstep{Local dimensions at the places above the coefficient prime.}
Let \(d=[L^+:\Q]\) again and let \(d_v\) be the local degree
at \(v\mid\ell\). Let \(d'_v\) be the rank of the pro-\(\ell\)
inertial units after removing their torsion. The local
abelianized quotient for the diagonal cocycles then has
rank \(d'_v+1\), including the unramified generator.
Here \(\sum_v d_v=\sum_v d'_v=d\), by local class field
theory on the chosen split places. Put
\[
 \nu=\binom{a}{2}+\binom{b}{2}+gab
\]
with \(a=n,b=0\) in type I. The primal first-order
coefficient module is \(B\) under conjugation by \(R\),
with \(-\tau\) at \(c\). Its rank is \(n+2\nu\).
At depth \(m\) the perfect pairing with the twisted dual is the composite 
\[
 (z,b)\longmapsto \langle z,b\rangle\longmapsto [t^{m-1}]\langle z,b\rangle\longmapsto \operatorname{Tr}_{K/\mathbf F_\ell}\!\left([t^{m-1}]\langle z,b\rangle\right).
\]

At \(v\mid\ell\), make the local odd coboundary changes
used in the choice of \(\kappa_j\) so that \(R\) is
upper triangular in the ordered local flag. The
local upper triangular coefficient subspace has
rank \(n+\nu\). The \(n\) diagonal cocycles kill the
local torsion prescribed by the level. Each of the
other positions, including the \(g\) copies for every
mixed-block position, is filtered by its distance
above the diagonal. For positions at a fixed diagonal distance, the higher-distance contributions vanish in the associated quotient. The action on the coordinate \((i,j)\) becomes
\[
 x_{ij}\longmapsto\psi_{v,i}(g)\,x_{ij}\,\psi_{v,j}(g)^{-1}\qquad(i<j).
\] The local
nonresonance test excludes the cyclotomic character
for these ratios. The finite module with this
character at depth \(m\) has order \(|K|^m\).

Tate local duality bounds its second cohomology by
\(|K|^{O(1)}\), since some ratio divided by the
cyclotomic character differs from one by a fixed
nonzero element of \(A\). The local Euler characteristic
formula and the formula for the size of coboundaries
then give \(|K|^{m(1+d_v)+O(1)}\) for its first
cocycles. Multiplying along the successive kernels gives
\[
       |Z_v|\le
       |K|^{m[n(d'_v+1)+\nu(1+d_v)]+C_v}.
\]
The image of \(Z_v\) in local cohomology quotients
by at least the coboundaries of the upper triangular
coefficient subspace. The upper triangular subspace has \(|K|^{m(n+\nu)}\) elements. Dividing by its coboundaries in the Greenberg--Wiles local term gives 
\[
 \frac{|L_v|}{|H_v^0|}\le \frac{|Z_v|}{|K|^{m(n+\nu)}}.
\]


\proofstep{The unramified and auxiliary factors.}
At each fixed non-\(\ell\) place of the unramified local
condition, the local term is one. At the new places
and at the distinguished test place the primal local
condition is full. Local Euler characteristic in
residue characteristic different from \(\ell\) makes
the local term the size of the local twisted-dual
invariants, and the dual local class must vanish.
At the auxiliary places the twist is trivial in
characteristic \(\ell\).

Write \(T\) for the regular Frobenius of the base curve. Its simple spectrum and the disjointness of the spectra in the two blocks give 
\[
 \dim_E\{X\in S:XT=TX\}=n,\qquad\{X\in M:XT=TX\}=0.
\]
 In the paired dual,
an invariant at the fraction field therefore has
\(\lambda_j=0\) for every \(j\) and at most \(n\)
free diagonal coordinates. For \(N\) sufficiently large as a function of \(m\), the finite-depth Frobenius equations agree with the fixed ones. The at most \(n\) free coordinates and the bounded Smith factors give the estimate
\[
 |K|^{mn}\,|K|^{C_j}=|K|^{mn+C_j} .
\]

\proofstep{Invariant spaces at the distinguished and real places.}
At the distinguished place the decomposition subgroup \(D_v\) must satisfy the intertwining test 
\[
 \bigl\{Y:\chi(h)r(h)Y=Yr(h)\text{ for every }h\in D_v\bigr\}=0.
\]
 This holds
over \(E\) by clearing denominators in \(Y\). Applying it
first to the off-diagonal projections and then to the
diagonal projection gives zero invariants for the paired
dual over \(E\), and a bounded finite-depth term.
For primal global invariants, commutation
with the span \(B\) forces centrality, while the
real action requires skew symmetry. The center consists of block scalars fixed by transpose. Since \(\ell\ne2\), the real skew condition gives 
\[
 Z\in Z(B),\quad Z=-Z^{\mathsf T}\quad\Longrightarrow\quad 2Z=0\quad\Longrightarrow\quad Z=0.
\]
 Hence the integral comparison contributes only bounded torsion at finite depth. At every real place
the primal invariant module has at least
\(\nu\) independent coordinates at each depth.
Indeed if \(c'\) is its conjugation, the matrix
\(R(c'c)\) is symmetric by the polarization
relation and \(\epsilon(c'c)=1\). Multiplication
by this matrix embeds the skew matrices in the
\(c'\)-fixed space.

\proofstep{The Greenberg--Wiles product.}
Greenberg--Wiles duality, with the tested dual bound,
gives the estimate in cohomology. The exponent of
\(|K|\) that multiplies \(m\) is
\[
  \sum_{v\mid\ell}
    \bigl(n(d'_v+1)+\nu(1+d_v)-(n+\nu)\bigr)
     +nq-\nu d=nd+nq .
\]
The primal cocycles normalized at \(c\) modulo their
normalized coboundaries give the same cohomology,
again by division by two. Thus the bound is
\(|K|^{m(nd+nq)+C}\) for all \(m\) and all sufficiently
large \(N\) at that \(m\). This estimate holds on the same auxiliary systems for every point in the chain.

\end{proof}

\subsection{The odd obstruction bound}
\label{subsec:odd-obstruction}

The first-order estimate counts directions. To bound the relations among them we must
also control the obstruction to lifting; it is the odd part that enters the two-block
calculation.

\begin{proposition}[Uniform control of the odd obstruction module]
\label{prop:marked-odd-obstruction-bound}
Use the fixed normalized point of Lemma~\ref{hyp:marked-square-zero-data}.
The module \(U_m\) is relative second cohomology of the odd matrix
module over \(K[[t]]/(t^m)\), with the local projections by the upper
triangular part at \(v\mid\ell\) and the unramified local condition
away from \(\ell\). A cocycle includes a trivialization of each
projected local class. There is a constant \(C\) independent of \(e,m\)
such that for every \(e\le m\)
\[
                  |U_m/t^e U_m|\le |K|^{ge+C}.
\]
In particular the quotient by the first \(g\) elementary-divisor
summands is killed by \(t^C\). Among any \(g+1\) classes there is
a combination with a unit coefficient whose multiple by \(t^C\)
vanishes.
\end{proposition}

\begin{proof}
\proofstep{The obstruction module at a fixed depth.}
Keep the fixed point and let \(A_e=A/(t^e)\).
Use the odd matrix module \(M_e\) over \(A_e\),
with conjugation by \(r\) and with \(-\mathsf T\)
at \(c\). At a place \(v\mid\ell\) the local
coefficient quotient \(Q_{v,e}\) is by the
upper triangular part in the ordered flag. Its
rank over \(A_e\) is \(ab\), one lower entry for each
pair of indices from different blocks. At the other
finite places the quotient is \(M_e\) itself and
the local condition is unramified.

In particular at a place above \(\ell\) the localization condition on the projected class \(\bar z\) is
\[
 \operatorname{res}_v[\bar z]=0\qquad(v\mid\ell).
\] Let \(S_e\)
denote the resulting first cohomology group with
these conditions. Let \(U_e\) denote second
cohomology with these local conditions.
Represent a class of \(U_e\) by a global
two-cocycle and a local one-cochain trivializing
its projection at every place. One quotients by the
changes from global one-cochains and local
one-cocycles with the allowed local restrictions.
Let \(r_e^2\) denote restriction followed by projection
in second cohomology, and set
\[
 a_e\colon H^1(G^+,M_e)\longrightarrow
             \bigoplus_v H^1(G_v,Q_{v,e})/L_{v,e}.
\]
Then the obstruction group sits in
\[
 0\longrightarrow\operatorname{coker}a_e\longrightarrow U_e
 \longrightarrow\ker r^2_e\longrightarrow0.
 \tag{11}\label{eq:local-patch-11}
\]

In \eqref{eq:local-patch-11} the local indices are precisely the primes above \(\ell\)
and the fixed tame primes dividing the level cutoff away from \(\ell\).
The distinguished place and the new Taylor--Wiles places belong to the
global ramification set. Only the new inertia evaluations enter \eqref{eq:local-patch-12}
in comparing with the base. Here \(G^+\) denotes the real-field group
\(G_{L^+,S'}\), the Galois group of the maximal extension unramified
outside the level set \(S'\) and the archimedean places. The subgroup \(L_{v,e}\) is the group unramified on the prescribed tame inertia kernel away from
\(\ell\) and zero above \(\ell\). The sequence follows
by first sending the cocycle to its global class.
If that class is zero one subtracts a global
coboundary and is left with the indicated local
classes modulo global restrictions.

At the base level,
\[
                         |S_e|\le |K|^{ge+C_0}.
 \tag{12}\label{eq:local-patch-12}
\]
Use evaluations on finitely many generators as for
the paired dual. For each local condition add the
coordinates of a local boundary matrix at the
places above \(\ell\), and require zero restrictions
on the prescribed inertias away from \(\ell\), where
the odd module is unramified. One obtains a
finite matrix of linear equations by taking the
row span, including all relations and local
equations. If \(F\) denotes this matrix, its nonzero Smith factors contribute a bounded correction: 
\[
 \log_{|K|}|\ker(F\bmod t^e)|=e\dim_E\ker(F\otimes_A E)+O_F(1).
\]
 Project to the generator
coordinates and divide by global boundaries.
The kernel of projection is given by the local
invariance equations on the added matrices,
so its count also has the corresponding
dimension estimate. This gives \eqref{eq:local-patch-12}, since the
resulting fraction-field quotient has dimension
\(g\).

At the auxiliary level, evaluate each class on
the new tame generators. Boundaries vanish there.
The Frobenius relation, the trivial inertia action
and \(Q_{j,N}\equiv1\pmod\ell\) make the
tame value fixed by conjugation under the
auxiliary Frobenius. Fix \(m\) and take \(N\) large enough for that depth. For every \(e\le m\) the tame value \(Z\) then obeys the same commutation equation as at the test element: 
\[
 Zr(h_j)=r(h_j)Z\pmod {t^e}.
\]
 The commutant of
\(r(h_j)\) has no off-diagonal part over \(E\)
because the full eigenvalue list is distinct.
By Smith normal form the cardinality of this evaluation image satisfies the bound
\[
 |\operatorname{im}({\rm ev})|\le |K|^{C_{\rm aux}}\qquad(e\le m).
\] In the kernel
the cocycle is zero on the new inertia and
descends to the base, with its old local
conditions. This proves \eqref{eq:local-patch-12} at the auxiliary
level with the additional constant \(C_{\rm aux}\).

Two more invariant groups have uniformly bounded
order. For \(H^0(G_v,Q_{v,e})\) above \(\ell\),
filter the lower entries so that the induced
characters are the lower ordered ratios. On a lower-entry graded piece, an invariant \(x_{ij}\) satisfies 
\[
 (\psi_{v,i}(h)\psi_{v,j}(h)^{-1}-1)x_{ij}=0\quad\text{for every }h.
\]
 Distinctness of the ordinary characters forces \(x_{ij}=0\) over the fraction field. Successive application to the filtration kills every invariant. The
same finite-row and Smith argument bounds
the finite-depth invariants. For the
cyclotomic twisted dual of \(M_e\), its invariants
on \(H\) vanish over \(E\). The twist is trivial
there, and conjugation invariants on the odd
matrices vanish by the block idempotents in
the span of \(r(H)\). The matrix trace pairing
identifies the odd module with its \(A_e\)-dual.

For \(X,Y\) in the odd matrix module the induced finite-field
pairing is
\[
 (X,Y)\longmapsto
 \operatorname{Tr}_{K/\mathbf F_\ell}\!\left([t^{e-1}]
                    \operatorname{tr}(XY)\right).
\]
Indeed a coefficient at degree \(j<e\) in a nonzero trace pairing
is detected by multiplication of the other vector by \(t^{e-1-j}\)
and a scalar of \(K\), using nondegeneracy of the field trace.
It gives a fixed
bound \(|K|^{C_d}\) for the global Tate dual
invariants. This remains valid at all auxiliary
levels because the coefficient action is
pulled back by a surjection. In particular
\[
          |\operatorname{coker} r^2_e|\le |K|^{C_d}.
 \tag{13}\label{eq:local-patch-13}
\]
Indeed take first the arithmetic level defining the global and local
groups. The finite set \(S'\) contains every place over \(\ell\) and
\(\ell\) is odd.

For the finite discrete \(\ell\)-primary coefficients
here global \(H^3\) vanishes, even in the presence of real places.
Local \(H^3\) at the finite places vanishes as well. Thus for any short exact sequence of these finite discrete coefficient modules, the end of the cohomology sequence is
\[
 H^2(G,A)\longrightarrow H^2(G,B)\longrightarrow H^2(G,C)\longrightarrow H^3(G,A)=0
\]
for \(0\to A\to B\to C\to0\). This is the required right exactness.
Poitou--Tate bounds the cokernel of
global localization with full local coefficients
by the global Tate dual invariants. Drop the components outside the fixed indices and identify each retained local group with the decomposition group at the chosen split prime of the CM field. For any retained local quotient \(T\twoheadrightarrow T'\) with kernel \(T_0\), local cohomological dimension two gives the exact segment 
\[
 H^2(D_v,T)\longrightarrow H^2(D_v,T')\longrightarrow H^3(D_v,T_0)=0.
\]

The level map from the arithmetic group at \(V_N\) to the group
for this count is a continuous bijection by the level-map
surjectivity and the injectivity in the arithmetic group comparison. Compactness and Hausdorffness make its inverse
continuous. Pullback of cochains thus identifies the global cohomology
and commutes with the fixed local restrictions, the signed real
involutions and the cyclotomic twist.  They precede passage to
the relative pro-\(\ell\) quotients for the patch.

\proofstep{The global and local Euler factors.}
\proofstep{Bounding the auxiliary odd classes by duality.}
The global Euler characteristic gives
\[
 |H^1(G^+,M_e)|=
    |H^0(G^+,M_e)|\,|H^2(G^+,M_e)|\,|K|^{eabd}.
\]
Each real conjugation has \(ab\) fixed coordinates
on \(M_e\). In the normalized basis its action
is minus transpose. At another real conjugation
multiplication by the corresponding invertible
symmetric block-diagonal matrix identifies the
fixed space with the skew off-diagonal matrices.
The local Euler term at \(v\mid\ell\) for
\(Q_{v,e}\) is \(|K|^{eab d_v}\) in the numerator
of \(|H^1|/(|H^0||H^2|)\). At the other
places \(|H^1_{\rm ur}|=|H^0|\), by taking the prime-to-\(\ell\) inertia invariants
and evaluating on Frobenius, as explained below. Taking orders
in \eqref{eq:local-patch-11}, using \(\sum d_v=d\) and dropping
the global \(H^0\) denominator, gives
\[
 |U_e|\le |S_e|\,
    \prod_{v\mid\ell}|H^0(G_v,Q_{v,e})|\,
                   |\operatorname{coker} r^2_e|.
 \tag{14}\label{eq:local-patch-14}
\]
In particular the two Euler contributions cancel
even if the finite-depth local quotients have
nonzero invariants.

To control elementary divisors, use both coefficient
sequences for \(e\le m\),
\[
 0\to A_{m-e}\xrightarrow{t^e} A_m\to A_e\to0,
 \qquad
 0\to A_e\xrightarrow{t^{m-e}} A_m\to A_{m-e}\to0.
\]
They stay exact on the odd modules and their local
quotients, since the local upper triangular kernel
has a free complement in its local frame. Write
\(i:U_{m-e}\to U_m\) for the first inclusion map
and \(r:U_m\to U_{m-e}\) for the second reduction
map. Multiplication and the first reduction satisfy 
\[
 i\circ r=t^e\operatorname{id}_{U_m},\qquad \ker(U_m\longrightarrow U_e)\subseteq i(U_{m-e}).
\]


To see this, subtract a lift of the coboundary
expression after reduction, global and local.
At a fixed tame place let \(D_u,I_u\) be its absolute
decomposition and inertia groups, and
\(J_u=I_u\cap\bigcap_i\ker\bar\chi_i\) for the residual diagonal
characters pulled back to \(D_u\). Its image on the base level is the
inertia kernel prescribed for the point. Thus \(J_u\) acts trivially on
the odd modules at every depth. The inertia quotient embeds into a product of multiplicative groups of the finite residual field \(k\). Therefore, for the number \(a\) of factors in the embedding,
\[
 |P_u|\mid (|k|-1)^a,\qquad (|P_u|,\ell)=1 .
\] A cocycle zero on \(J_u\) descends to
\(D_u/J_u\). Its restriction to \(P_u\) is a boundary by averaging.

Subtract that boundary and obtain a cocycle on
\(D_u/I_u\simeq\widehat{\mathbb Z}\) with values in the
\(P_u\)-invariants. Invariants lift under a surjection by the same
averaging. Lift the value on Frobenius, which gives a continuous
cocycle since the cyclic subgroup generated in the finite affine
group is finite. Add back a lift of the boundary to lift the original
one-cocycle exactly. This calculation also identifies the local
\(H^1\) unramified on \(J_u\) with Frobenius coinvariants on
the \(P_u\)-invariants. On the finite invariant coefficient module \(M\), let \(T=\mathrm{Frob}-1\).
Rank counting for this endomorphism gives
\[
 |\operatorname{coker}T|=\frac{|M|}{|\operatorname{im}T|}
                         =|\ker T|.
\]
Thus the Frobenius coinvariants and invariants have the same order. At the
\(\ell\)-adic places the local allowed classes
are boundaries and their boundary matrices
lift. The remaining cochains are then in the
coefficient kernel and hence give the inverse
image under \(i\).

For the second sequence,
\(|\operatorname{coker}r|\) is bounded by
\(|\operatorname{coker}r^2_e|\). Lift the global class
using surjectivity on global second cohomology.
The defects of its local nullhomotopies now have
coefficients in \(A_e\). For the finite, discrete coefficient modules a section of the coefficient surjection is continuous as a map of sets. A local one-cochain \(b\) therefore has a continuous lift by composition:
\[
 \tilde b=s\circ b,\qquad \operatorname{pr}\circ\tilde b=b .
\]
The defects are second cohomology classes in
the local quotients at depth \(e\). Changing the
global lift changes them by the image of
\(r^2_e\). If the defect in its cokernel vanishes,
adjust the global lift accordingly and use the
local cochains of the corrections; this produces
a lift in \(U_m\). The same comparison on
differences shows that equal defects yield the
same element of \(\operatorname{coker}r\).
This is the local restriction cokernel bound,
using global cohomological dimension two.

\proofstep{The quotient at a smaller depth and the elementary divisors.}
Consequently
\[
 |U_m/t^eU_m|\le |U_e|\,|\operatorname{coker} r^2_e|
                  \le |K|^{ge+C},\qquad e\le m.
 \tag{15}\label{eq:local-patch-15}
\]
Here \(C\) is the sum of the constants in \eqref{eq:local-patch-12},
the auxiliary correction, the local invariant
bounds and twice the constant in \eqref{eq:local-patch-13}. It is
chosen before \(m\) or \(e\). Write the elementary
divisors of \(U_m\) as \(A/(t^{h_i})\) in
nonincreasing order, with \(h_i\le m\).
If there is a \((g+1)\)-st divisor, use its own exponent as the test depth in \eqref{eq:local-patch-15}:
\[
 e=h_{g+1}\quad\Longrightarrow\quad h_{g+1}\le C .
\]
Thus the quotient by the first \(g\) summands is
killed by \(t^C\).

For any \(g+1\) obstruction
classes, their projections to the first \(g\)
summands admit a linear relation with at least
one unit coefficient. For example lift to a
matrix of \(g\) rows over \(A\), take a
fraction-field linear dependence and rescale to
integral coefficients with a unit. Multiplying
this combination by \(t^C\) kills the full
obstruction. For \(g=0\) \eqref{eq:local-patch-15} directly gives the
same annihilation bound.  Its obstruction classes
are the odd two-cocycles with their local
trivializations. Their coboundaries are precisely
the global one-cochain changes with the same
local conditions, as in \eqref{eq:local-patch-11}.

\end{proof}

\subsection{One marked system for the chain}
\label{sec:patching-marking}
\hypertarget{source-ordinary-patching-marking}{}

We now put the bounds at the individual points to work on a single system. Keep the
ordinary determinant and the finite chain of comparison arcs fixed. The group \(\Gamma_0\) has a positive subgroup cut out by
the CM sign. The residual characters are defined on the positive
subgroup.
In the CM subgroup write \(\bar\chi_i\) for the residual characters just specified. The subgroup used in the marking is the intersection
\[
 H_{\rm res}=\ker\bar\epsilon\cap\bigcap_i\ker\bar\chi_i
\] It is open.
Polarization permutes the inverse residual characters up to their
common cyclotomic multiplier, so this joint kernel is normal.
The finite tests of Section~\ref{subsec:det-test}
and the simple Frobenius spectra are imposed at all the arcs
before choosing the auxiliary levels. The resulting marked
Frobenius elements \(\sigma_j\) lie in the cyclotomic kernel.

At level \(N\), the surjection \(p_N:\Gamma_N\to\Gamma_0\)
identifies the fixed decomposition maps and the real involution.
It takes the tame generator \(\tau_{j,N}\) to \(1\)
and sends the marked Frobenius to \(\sigma_j\) modulo \(H_{\rm res}\)
and arbitrarily close to \(\sigma_j\) on finite quotients of
the pro-\(\ell\) system. The inertia generator and the norm consequently satisfy 
\[
 p_N(\tau_{j,N})=1,\qquad Q_{j,N}\equiv1\pmod{\ell^N}.
\]
 These requirements on split primes are
obtained by Chebotarev in the finite quotients specifying
the Frobenius approximations, including the cyclotomic quotient
of depth \(N\). The element being approximated is already
trivial in that cyclotomic quotient.

\begin{proposition}[A common marked profinite limit for the chain]
\label{prop:marked-profinite-system}
Fix the comparison arcs and auxiliary prime systems of
Proposition~\ref{prop:marked-uniform-dual-tests}, with each
auxiliary Frobenius converging to its test element in the
common quotient and each norm congruent to \(1\bmod\ell^N\).
After passing to the common residual-cut pro-\(\ell\) quotients,
the point actions, ordinary determinants and representations giving
nonsemisimple quotients at thick points factor through the reduced marked groups. There are a finitely
generated profinite group \(G\) and a surjection
\(p:G\twoheadrightarrow\Gamma_0\) compatible with the fixed
local maps and signs. The distinguished pairs satisfy
\[
 p(\phi_j)=\sigma_j,\qquad
 p(\tau_j)=1,\qquad [\phi_j,\tau_j]=1.
\]
The kernel is the closed normal subgroup generated by the \(\tau_j\).
All these comparisons can be read simultaneously at any fixed finite
depth, with the auxiliary levels allowed to depend on the depth.
\end{proposition}

\begin{proof}
\proofstep{Killing the new tame inertia in the residual cut.}
For a profinite group \(J\) write \(P_\ell(J)\) for the
kernel of its maximal pro-\(\ell\) quotient. Put
\(H_N=p_N^{-1}H_{\rm res}\). The kernel of \(p_N\) is contained
in the closed normal subgroup generated by the
\(\tau_{j,N}\) and \(P_\ell(H_N)\). In the Galois description,
killing the new inertia gives the original ramification
conditions. The additional pro-\(\ell\) quotient of the
pullback of \(H_{\rm res}\) retains just the pro-\(\ell\) tame generator
at each new place.

At residue characteristic \(p\ne\ell\), any image of a pro-\(p\) group in a pro-\(\ell\) quotient is trivial, as its finite quotient orders would have to be powers of both primes:
\[
 p^a=\ell^b\quad\Longrightarrow\quad a=b=0 .
\]
This kills wild inertia and, for the same reason, every prime-to-\(\ell\) primary component of tame inertia. Each new split place and its
conjugate give the same normal generator condition.
The local tame relation is
\(\phi_{j,N}^{-1}\tau_{j,N}\phi_{j,N}=\tau_{j,N}^{Q_{j,N}}\)
in this quotient. Here conjugation by \(\phi_{j,N}^{-1}\) acts on the tame generator by the residue cardinality \(Q_{j,N}\).

\proofstep{The pro-\(\ell\) descent at each depth.}
At every comparison arc the characteristic polynomial on \(H_{\rm res}\)
reduces to \((X-1)^n\). The residual image of \(H_{\rm res}\) in the
finite field matrix group consists of unipotents. If \(u=1+N\) in this finite image, then \(N^n=0\) and in
residue characteristic \(\ell\)
\[
 u^{\ell^a}=1+N^{\ell^a}=1\qquad(\ell^a\ge n).
\]
Every element has \(\ell\)-power order. Cauchy's theorem then makes the
whole image an \(\ell\)-group. Each subsequent congruence kernel in
the curve representation is an \(\ell\)-group. Thus the point
action kills \(P_\ell(H_{\rm res})\), with the same \(H_{\rm res}\) for all points.
The same argument applies to the paired and affine actions at a
finite truncation, using the additive kernels of the coefficient
maps.

For continuous determinants at the ordinary levels,
the split multiplicity-free residual characters give the same
factorization. One may see it on the finite Cayley--Hamilton
quotients: writing \(\bar h\) for the group-algebra image of \(h\in H_{\rm res}\) modulo the maximal ideal, its characteristic polynomial yields 
\[
 P_h(X)\bmod\mathfrak m=(X-1)^n,\qquad (\bar h-1)^n=0.
\]
 At each successively deeper finite quotient the image
of \(H_{\rm res}\) therefore has \(\ell\)-power order. The determinant coefficients and ordered local characters thus factor through the respective quotients
\[
 \Gamma_0\twoheadrightarrow\Gamma_0/P_\ell(H_{\rm res}),\qquad \Gamma_N\twoheadrightarrow\Gamma_N/P_\ell(H_N).
\]
\label{cor:ordinary-patching-marking-1}

\proofstep{The closed matrix algebras after descent.}
This reduction also preserves the thick links. The image of every
element of \(P_\ell(H_{\rm res})\) in the integral closed
Cayley--Hamilton algebra is one at every coefficient depth, hence
one in the separated algebra. Its continuous group map factors
through the quotient group. This gives a map from the closed
Cayley--Hamilton algebra of the quotient group back to the original
one: the determinant and local ordered relations descend, and the
kernel of the map to the finite integral algebra is closed.

In the
opposite direction the quotient homomorphism of groups carries
each defining relation to its counterpart and respects closure.
Writing \(\mathcal E(\Gamma)\) for this closed algebra, the
two constructions give
\[
 \mathcal E(\Gamma_0)\ \longrightarrow\
 \mathcal E(\Gamma_0/P_\ell(H_{\rm res}))\
 \longrightarrow\ \mathcal E(\Gamma_0).
\]
The compositions fix the coefficients and each group image.
These topologically generate the algebras, so both compositions
are the identity. Extending scalars to the field of a nonsemisimple quotient preserves that quotient. The argument
also applies to the marked group with \(H_N\) in place of
\(H_{\rm res}\).

\proofstep{A uniform Frattini bound.}
The bound on the number of generators for the pro-\(\ell\)
groups \(H_N/P_\ell(H_N)\) is uniform in \(N\).
In this Galois setting the fixed field of the residual cut has
fixed degree. Class field theory bounds the \(\mathbb F_\ell\)
dimension of the unramified characters by its fixed class group,
and the restrictions at the fixed finite places contribute a
fixed finite dimension by local class field theory. At each new tame place \(w\), the pro-\(\ell\) tame inertia is procyclic. Its contribution is bounded by 
\[
 \dim_{\mathbb F_\ell}\operatorname{Hom}_{\mathrm{cont}}(I_w,\mathbb F_\ell)\le1.
\]

Each of the at most \(2q\) new places of the CM base has at most the fixed extension degree \(d'\) places above it. Hence
\[
 \#\{\text{new places upstairs}\}\le d'\,\#\{\text{new places downstairs}\}\le 2d'q .
\]
The kernel of restriction to all allowed inertia is unramified.
The same bound with no new places applies to \(H_{\rm res}/P_\ell(H_{\rm res})\).
For a pro-\(\ell\) group \(P\), the continuous characters detect exactly its Frattini quotient: 
\[
 \operatorname{Hom}_{\rm cont}(P,\mathbb F_\ell)=\operatorname{Hom}_{\rm cont}\!\left(P/\overline{P^\ell[P,P]},\mathbb F_\ell\right).
\]
 From now on use the quotients, so \(H_N\) and \(H_{\rm res}\)
themselves are pro-\(\ell\) and finitely generated.

\proofstep{One group marking every auxiliary level.}
Choose a free profinite group \(\mathcal F\) with one generator
for each element of the finite quotient \(\Gamma_0/H_{\rm res}\), a
separate generator for the real involution, and enough additional
generators for the largest of these bounds on \(H_N\). Send
the first generators to representatives in \(\Gamma_N\), the
involution generator to the fixed involution, and the remaining
generators to generators of \(H_N\), padding with identities.
These choices give surjections \(\pi_N:\mathcal F\to\Gamma_N\)
with one common composition into \(\Gamma_0/H_{\rm res}\).
Let \(\mathcal U\) be an ultrafilter extending the tails of the
positive integers.

\proofstep{Compatible finite-quotient limits.}
At an open normal subgroup \(V\subset\mathcal F\), only finitely many subgroups of \(\mathcal F/V\) occur as images. The ultrafilter selects the unique \(K_V\) satisfying
\[
 \{N:(\ker\pi_N)V/V=K_V\}\in\mathcal U .
\]
They are compatible, and every projection between two such
subgroups is onto because it is onto for the corresponding
kernels on their common ultrafilter event. Consequently
\[
 K=\varprojlim_V K_V,\qquad G=\mathcal F/K,\qquad
 G/\operatorname{im}(V)=\mathcal F/(KV).
\]
Compactness and surjectivity of the transition maps give
\(K\bmod V=K_V\); closed normality follows from that of the
\(K_V\).

For every finite quotient \(Q\) of \(\Gamma_0\), the
maps
\[
 \mathcal F\xrightarrow{\pi_N}\Gamma_N \xrightarrow{p_N}\Gamma_0\to Q
\]
have finitely many
possibilities. Their ultrafilter value kills \(K\), since at
the corresponding finite depth the images of the level kernels
coincide with the image of \(K\). The resulting maps \(G\to Q\)
are compatible, so define the continuous map \(p:G\to\Gamma_0\). The map is onto all finite quotients, so it is onto
by compactness. The CM sign, cyclotomic character and real
involution on \(G\) are thereby compatible with the base and
have the required finite quotient comparisons.

\proofstep{Local maps and the marked tame relation.}
The fixed local restriction maps also have simultaneous limits.
For the fixed local group \(D\), let \(f:D\to\Gamma_0\) be its given map. Its subgroup over the residual kernel is open, since
\[
 [D:f^{-1}(H_{\rm res})]=[f(D):f(D)\cap H_{\rm res}]\le[\Gamma_0:H_{\rm res}]<\infty .
\] The maps to all the
\(\Gamma_N\) kill its pro-\(\ell\) kernel since it maps into
\(H_N\). The resulting quotient of the local group is finitely
generated. Its open pro-\(\ell\) subgroup has finite Frattini
rank by local class field theory, and the complementary index is
finite. At each fixed \(V\), equality of the kernel images on an ultrafilter event gives 
\[
 \Gamma_N/\pi_N(V)\simeq\mathcal F/(V\ker\pi_N)=\mathcal F/(VK).
\]
 The finite set of maps from the finitely
generated local quotient to this target has one ultrafilter
value.

Taking all \(V\) defines the local map to \(G\)
with agreement for every element of the local group on a
single event at each fixed depth. Apply the same argument to
the finite tuple of elements \((\phi_{j,N},\tau_{j,N})\)
to obtain \((\phi_j,\tau_j)\) in \(G\). The markings at the limit have the specified Frobenius and inertia values
\[
 p(\phi_j)=\sigma_j,\qquad p(\tau_j)=1 .
\] The \(\tau_j\) lie in \(p^{-1}H_{\rm res}\), a
pro-\(\ell\) group because its finite images come from the
\(H_N\). In a fixed finite quotient choose \(h\) so that the order of the image \(\bar\tau_j\) divides \(\ell^h\). For \(N\ge h\) the finite-depth tame relation
passes to
\[
 \bar\phi_j^{-1}\bar\tau_j\bar\phi_j
 =\bar\tau_j^{\,Q_{j,N}}=\bar\tau_j
 \qquad (Q_{j,N}\equiv1\bmod\ell^h).
\]
The commutator is trivial in every finite quotient and hence in \(G\).

\proofstep{Recovering the base quotient from the limit.}
Finally, every finite quotient of \(G\) killing all the
\(\tau_j\) descends to \(\Gamma_0\). Indeed read its quotient
map on \(\mathcal F\) at a fixed finite depth. There it factors
through \(\Gamma_N\) and kills each marked \(\tau_{j,N}\)
on an ultrafilter event, hence factors through \(p_N\).
The maps from the finitely generated group \(\Gamma_0\)
to this finite quotient form a finite set. Its selected
ultrafilter value therefore gives the factorization through
\(p\).

Put \(N_\tau=\overline{\langle\!\langle\tau_j
\rangle\!\rangle}\). The markings give \(N_\tau\subset\ker p\).
The factorization shows
\[
 \ker p\ \subset\
       \bigcap_{\substack{Q\ {\rm finite}\\G/N_\tau\twoheadrightarrow Q}}
       \ker(G\longrightarrow Q)=N_\tau ,
\]
where the equality holds because \(G/N_\tau\) is profinite.
Thus \(\ker p=N_\tau\). This gives the descent used for
the odd cocycles in Section~\ref{sec:patching-cotangent}.
The construction uses the same free group, limits and local
maps at all chain points. The ultrafilter matching can depend
on truncation depth, and no single finite auxiliary level is
required to match all depths.
\label{cor:ordinary-patching-marking-2}
\end{proof}

\subsection{The polarized matrix algebra at a comparison arc}
\label{sec:patching-matrices}
\hypertarget{source-ordinary-patching-matrices}{}

The marking keeps the tests fixed along the chain. To recover matrix entries from
determinant data, begin with the common patch ring. It carries the universal group determinant on the
free marked cover, the ordered local characters, the ordered
Frobenius roots, the tame variables and one coefficient variable.
In its quotient \(B_0\) impose the kernel of the marking, polarization,
ordinary and inertia identities and the auxiliary trace identities against
all group elements. Its augmentation maps onto the ordered splitting
algebra of the original determinant ring. Indeed the images contain the
original determinant coefficients, ordered characters, coefficient
variable and the adjoined ordered roots, which generate that algebra.
The group and character maps give the augmentation, and each auxiliary
identity maps to the Cayley--Hamilton identity for an actual split
characteristic polynomial or to zero on tame inertia.

\begin{lemma}[Common coefficients on a marked chain]\label{lem:common-marked-coefficients}
A finite marked chain of comparison arcs admits one enlarged
finite residue field with all ordered roots and link
inclusions preserved. The arc maps are surjections to
\(k[[t]]\). At an arc with comparison ring \(B\), the marked
coefficient map extends to the completion of
\(O[[S]]_{(\varpi)}\) and identifies its residue field by
\[
 \operatorname{Frac}(k[[S]])\xrightarrow[\ S\mapsto t\ ]{\ \sim\ } k((t)).
\]
With this identification the residual coefficient map is the identity.
\end{lemma}

\begin{proof}
\proofstep{Coefficient extensions at a marked arc.}
Here are the coefficient choices implicit in completing at a point.
Enlarge the residue fields of the finite list of arcs to one finite field
\(k\), extend the coefficient DVR by the corresponding unramified extension
and adjoin \(s\), evaluated at a uniformizer \(t\) of each arc.
Write \(Q_1\supseteq P\subseteq Q_2\) for a link between adjacent curve kernels. In the finite flat coefficient extension start with a lift \(Q'_1\) of \(Q_1\). Going down and then going up give primes with contractions 
\[
 Q'_1\supseteq P'\subseteq Q'_2,\qquad P'\cap R=P,\quad Q'_2\cap R=Q_2,
\]
 where \(R\) denotes the ring before coefficient extension. Lying over applies also to a chain with just
one point.

The monic polynomial lifting the separable residue extension splits over each arc
after the chosen enlargement, so the specified prime over the curve is
the kernel of one of the evaluations. Extending each link by power series
in \(s\) retains a domain quotient and its inclusion in both adjacent
evaluation kernels. Apply the same zigzag lifting to the finite free
ordered splitting algebra. For a monic \(P(T)\) of degree \(n\), division with remainder over \(S\) gives the free module decomposition for the first root algebra
\[
 S_1=S[T_1]/(P(T_1))=\bigoplus_{j=0}^{n-1} ST_1^j .
\] In \(S_1[T]\) one has
\[
 P(T)=(T-T_1)P_1(T),\qquad \deg P_1=n-1 .
\]
Repeating with \(P_1\) gives the ordered splitting algebra, of rank \(n(n-1)\cdots1\) before selecting the prescribed factor. Select the local direct factor
with the prescribed residual ordering by lifting its residue
idempotent.

It remains finite free over the complete local ring.
Here each actual characteristic polynomial
splits with distinct roots in the fraction field of the arc, and those
roots lie in its DVR by integrality. Tensoring the fiber at the generic
point of the arc with that field gives a product indexed by the compatible orderings.
Thus every chosen prime over the arc picks an evaluation with those
roots. The maps to the arc contain the constants and send \(s\) to \(t\). Completeness supplies every series by
\[
 \sum_{i\ge0}a_i s^i\longmapsto \sum_{i\ge0}a_i t^i\qquad(a_i\in k),
\]
so the maps are surjective. Inclusion of the original Hecke ideal can be tested after
contraction. At a link with nonsemisimple residue algebra,
the field extension and the change of
group require the argument in Section~\ref{subsec:closed-link}.

Fix one such arc and let \(a:B_0\twoheadrightarrow k[[t]]\)
be the comparison map. Write \(\mathfrak p=\ker a\),
\(B=\widehat{(B_0)_{\mathfrak p}}\) and \(F=k((t))\).
The residue map \(B\to F\) extends \(a\).
Let \(O\) be the coefficient DVR with uniformizer \(\varpi\),
surjecting onto \(k\), and let \(s\in B_0\) evaluate to \(t\).
\proofstep{The coefficient DVR at the comparison point.}
The continuous substitution \(O[[S]]\to B_0\) sends \(S\) to
\(s\). Composing with \(a\) reduces the coefficients and
sends \(S\) to \(t\), so the kernel of the composition is
\((\varpi)\). For \(D_0=O[[S]]_{(\varpi)}\), this substitution gives
\[
 D_0/\varpi D_0=\operatorname{Frac}(k[[S]])
        \xrightarrow[\ S\mapsto t\ ]{\ \sim\ } k((t))=F .
\]
Every denominator outside \((\varpi)\) thus maps to a unit of \(B\); the map factors through \(D_0\). Since \(\varpi^h D_0\) maps into \(\mathfrak m_B^h\) for all \(h\) and \(B\) is complete, it extends to \(\mathcal O=\widehat{D_0}\). The identification
of its residue field with \(F\) is induced by the same
substitution. Lifts of elements of
\(k[[t]]\) can be taken in \(O[[S]]\), so their images in \(B\)
also come from \(B_0\).
\label{cor:ordinary-patching-matrices-1}

\end{proof}

\paragraph{The finite Cayley--Hamilton quotient.}
To pass from traces to an algebra acting on a module, start on the CM subgroup \(G_{\rm CM}\) of Section~\ref{sec:patching-marking}
take the closed Cayley--Hamilton quotient \(E_0\) of
\(B_0[G_{\rm CM}]\). Take the closed ideal generated by the ordered ordinary, prescribed inertia and auxiliary relations together with their images under polarization. The determinant descends through that ideal by
the trace-kernel calculation below. The resulting algebra \(E_0\) carries a determinant and an anti-involution. On a group element the latter sends \[ g\ \longmapsto\ \epsilon(g)^{1-n}(cgc^{-1})^{-1}. \] Here \(c\) is the real involution from the marking. The polarization is expressed in the frame in which the anti-involution specializes to transpose at the arc point. We use degree-\(n\)
determinants as multiplicative polynomial laws, with their
Cayley--Hamilton identity under arbitrary scalar extension.
\cite{chenevier}\label{cor:ordinary-patching-matrices-2}

\begin{auxlemma}[Detecting the full matrix span by traces]\label{lem:matrix-span-by-traces}
Let \(k\subseteq k'\) be fields, \(n\ge1\), and let \(S\subseteq M_n(k')\) satisfy \(\operatorname{tr}(XY)\in k\) for all \(X,Y\in S\).
The \(k'\)-span of \(S\) equals \(M_n(k')\) if and only if some \(X_1,\ldots,X_{n^2}\in S\) satisfy
\[
 \det\bigl(\operatorname{tr}(X_aX_b)\bigr)_{1\le a,b\le n^2}\ne0\quad\text{in }k.
\]
\end{auxlemma}

\begin{proof}
The matrix units have trace-dual basis \((E_{ji})\), since
\[
 \operatorname{tr}(E_{ij}E_{uv})=\delta_{ju}\delta_{iv}.
\]
Thus a basis selected from a spanning set has nonsingular trace matrix. Conversely a nonsingular trace matrix forces the \(n^2\) selected elements to be independent over \(k'\), so they span. The determinant belongs to \(k\), and the inclusion into \(k'\) preserves nonvanishing.
\end{proof}

\begin{proposition}[The finite polarized Cayley--Hamilton algebra]
\label{prop:marked-finite-ch-algebra}
Over the completed arc comparison \(B_0\) built from
Proposition~\ref{prop:marked-profinite-system}, the closed
Cayley--Hamilton quotient \(E_0\) of \(B_0[G_{\rm CM}]\) is finite
with its adic topology, polarized anti-involution and determinant law. When the residue characters are distinct, lifts of their diagonal
idempotents can be chosen self-adjoint. In the two-block case let
\(M\) be the mixed corner module over the coefficient ring \(B\).
At every specialization with coefficient kernel \(Q\) admitting a nonsemisimple residue quotient one has
\[
          \operatorname{Ann}_B(M)\subseteq Q .
\]
\end{proposition}

\begin{proof}
\proofstep{Finiteness and the adic topology of the quotient.}
We record the finiteness and continuity needed for the matrix
realization. The original residual determinant on \(G_{\rm CM}\)
is the direct sum of \(n\) distinct characters over the finite
residue field of \(O\). In a finite group quotient through
which they factor, let \(C\) be the corresponding residual
Cayley--Hamilton algebra. It surjects onto the product of the
\(n\) scalar factors, since distinct characters are linearly
independent. Its kernel \(J\) consists of nilpotents by the
Cayley--Hamilton equation, so \(J\) is its radical in this
finite-dimensional algebra. The diagonal units lift to
orthogonal idempotents.

For a lifted rank-one idempotent the determinant law and its corner satisfy
\[
 D(1+Xe_i)=1+X,\qquad e_iCe_i=k_0 e_i,
\]
where \(k_0\) is the residual field. These are the idempotent rank
and corner rules for determinants. Expand a product in \(J^n\) by inserting \(\sum_i e_i=1\) between
its factors. A term has indices \(i_0,\ldots,i_n\). Some
\(i_a=i_b=i\) with \(a<b\), and the corresponding subproduct lies in
\[
 e_iJ^{\,b-a}e_i\subseteq e_iJe_i=0.
\]
Every summand vanishes, hence \(J^n=0\).
In \(J/J^2\) a linear functional on the \((i,j)\)-corner
produces a cocycle for \(\chi_i\chi_j^{-1}\) by reading the
\((i,j)\) entry of a group element and dividing by \(\chi_j\).

For the chosen functional write \(x(g)\) for its value on that
entry and \(z(g)=x(g)/\chi_j(g)\). Since \(J^2=0\) in the first-order
quotient, multiplication gives
\[
 x(gh)=\chi_i(g)x(h)+x(g)\chi_j(h),\qquad
 z(gh)=z(g)+(\chi_i\chi_j^{-1})(g)z(h).
\]
This map is injective on the dual corner, since the images
of the group elements span \(C\), and \(i\ne j\) on each
nonzero radical corner. The dimension of each space of
these continuous cocycles on \(G_{\rm CM}\) is bounded by \(h\),
uniformly in the finite quotient. For example their values
at a finite topological generating set determine them,
and \(G_{\rm CM}\) is open in a finitely generated profinite group.
Products of radical lifts span every \(J^j/J^{j+1}\).
It follows that
\[
                 \dim_{k_0} C\le n+\sum_{j=1}^{n-1}(n^2 h)^j .
\]

For each coefficient truncation, adding \(u-1\) for all \(u\)
in an open normal subgroup gives an algebra quotient of a
finite group ring. Its reduction modulo the coefficient
maximal ideal has the preceding bound once that subgroup
lies in the residual character kernels. Lifting spanning
group elements and using Nakayama gives the same number
of generators over each coefficient truncation. At a fixed depth \(b\) write \(C_{U,b}\) for the resulting quotient at an open normal subgroup \(U\). If at most \(d\) generators suffice over the finite scalar ring \(R_b\), then 
\[
 U'\subseteq U\ \Longrightarrow\ C_{U',b}\twoheadrightarrow C_{U,b},\qquad |C_{U,b}|\le |R_b|^d.
\]
 After a quotient attains the maximal size, every smaller subgroup gives an isomorphism.

To choose one
set of group elements, take a spanning set after stabilization
at residue depth. Every deeper group and coefficient quotient
with that residue quotient is generated by the same elements,
again by Nakayama. For each element of the closed quotient
its expressions with these generators at finite depth define
nonempty closed sets in a fixed finite power of \(B_0\),
which is compact. Let \(C_b\) be the closed set of coefficient vectors representing the chosen element at depth \(b\). Compatibility and compactness give
\[
 C_{b+1}\subseteq C_b,\quad C_b\ne\varnothing\ (b\ge1)\quad\Longrightarrow\quad \bigcap_{b\ge1} C_b\ne\varnothing .
\] This proves finite generation in the closed quotient.
At each coefficient depth the stabilized group quotient
also shows that \(u-1\) lies in that coefficient power times
\(E_0\) on an open normal subgroup.  Take the closed quotient using the finite group and coefficient quotients. \label{cor:ordinary-patching-matrices-3}

\proofstep{Checking the trace kernel at finite precision.}
For the trace-kernel calculation first take the finite
closed Cayley--Hamilton quotient without the extra
relations. Over \(B_0\), the split residual characters
give \(n\) rank-one idempotents with scalar corners,
by the idempotent lifting argument below applied to
the original residue field. The trace radical in this quotient \(C\) is the common kernel
\[
 \bigcap_{x\in C}\ker\bigl(z\longmapsto\operatorname{tr}(xz)\bigr).
\]
Cyclicity of trace makes it a two-sided ideal. Its diagonal corners are
zero by taking \(x=e_i\). These properties hold for the
image of the ideal after any scalar extension. The
determinant of \(1+z\) for a member of this ideal is
one: use elimination in the rank-one corners. Each pivot
is one because all corrections are in the trace radical.
Off-diagonal transvections have determinant one by scaling
their entry by a conjugation over a Laurent polynomial
ring and specializing the polynomial identity at zero.

For general \(x\), let \(P_x(X)=D(X+x)\), a monic
polynomial. Cayley--Hamilton makes \(X+x\) invertible after
localizing at \(P_x\), and the trace radical is an ideal.
Thus multiplicativity in that localization gives
\[
 D(X+x+z)=D(X+x)\,
                  D(1+(X+x)^{-1}z)=D(X+x).
\]
Multiplication by a monic polynomial is injective over
any scalar ring. The equality therefore holds before
localization, and \(X=0\) proves invariance on each
trace-radical coset as a polynomial law. Thus a relation whose trace pairings
with all group elements vanish is killed as a polynomial
law, in arbitrary residue characteristic.
 The ordered ordinary and fixed inertia
generators are the relations of the ordinary pair, whose
polarized images also lie in this ideal. The closure
preserves the relations by continuity. \label{cor:ordinary-patching-matrices-4}

Extend scalars to the completed point ring. Finiteness over \(B\) identifies the resulting algebra with its completion: 
\[
 E=B\otimes_{B_0}E_0\simeq\varprojlim_b E/\mathfrak m_B^bE.
\]
 The determinant and anti-involution extend by scalar extension. There is a homomorphism
\(r:E\to M_n(F)\) given by the specialized point
representation. The point kills the defining relations,
so this homomorphism extends the given action on \(G_{\rm CM}\).
Let \(n_a\) denote the constituent dimensions. The positive cyclotomic kernel already spans the full diagonal blocks at the arc. Since \(B\to F\) is onto, 
\[
 r(E)=\bigoplus_a M_{n_a}(F)\subseteq M_n(F).
\]

The determinant of \(r(x)\) and its characteristic
polynomial are the reductions of the corresponding
polynomial law for each \(x\in E\). For \(x\in\ker r\) its residual characteristic polynomial has only the term of highest degree, so Cayley--Hamilton gives
\[
 x^n\equiv0\pmod{\mathfrak m_B E}.
\]

\paragraph{Self-adjoint rank-one idempotents.}
We need idempotents compatible with the polarization. Lift the diagonal matrix units
successively, each time in the
corner complementary to the idempotents already lifted.
For a lift \(x\) of a diagonal unit, averaging \(x\) with
\(x^*\) makes it self-adjoint, since \(2\) is a unit.
The defect \(x^2-x\) lies in \(\ker r\) and is topologically
nilpotent by the previous calculation. Write \(d=x^2-x\). Since \((2x-1)^2=1+4d\), the inverse of
\(2x-1\) exists by completeness and commutes with \(x\).
One Newton step gives
\[
      x'=x-(2x-1)^{-1}d,\qquad (x')^2-x'=(2x-1)^{-2}d^2 .
\]
Iteration squares the defect up to a unit at every step.
Its limit is a self-adjoint idempotent of the same reduction.

In successive complementary corners the resulting lifts
are orthogonal, and the final complement itself lifts the
last diagonal unit. Denote the family by \(e_i\).
For an idempotent \(e_i\) the determinant on \(1+Xe_i\) is
\((1+X)^{r_i}\) over the local ring. The leading coefficient of \((1+X)^{r_i}\) is one, so reduction
preserves its degree. Since the reduction is \(1+X\), we have \(r_i=1\).
For \(y\in e_iEe_i\), its degree-one characteristic polynomial
in this corner is \(T-\operatorname{tr}(y)\). The corner unit is
\(e_i\), so Cayley--Hamilton reads
\[
 y-\operatorname{tr}(y)e_i=0,\qquad e_iEe_i=B e_i .
\] Its trace sends \(e_i\) to \(1\).
 \cite{chenevier}\label{cor:ordinary-patching-matrices-5}

\proofstep{Lifting and matching the matrix units.}
Choose a reference index \(o(i)\) for each block, constant
throughout that block. Lift the matrix unit \(E_{o(i),i}\)
to \(z_i\in e_{o(i)}Ee_i\). The two products belong to the rank-one corners at \(i\) and
\(o(i)\). Their coefficients coincide because
\[
 \operatorname{tr}(z_i^*z_i)=\operatorname{tr}(z_i z_i^*),
 \qquad \operatorname{tr}(e_i)=\operatorname{tr}(e_{o(i)})=1.
\]
Write \(b_i\) for this common coefficient. Its residue
is \(1\). Taking its unique square root congruent to \(1\)
and dividing \(z_i\) by that root gives \(u_i\) with
\[
            u_i^*u_i=e_i,\qquad u_i u_i^*=e_{o(i)} .
\]
For the reference index use \(u_i=e_i\).  In particular
\[
      x\longmapsto (u_i x u_j^*)_{ij},
      \qquad (x_{ij})\longmapsto\sum_{i,j}u_i^*x_{ij}u_j
\]
are inverse coordinate changes. Each \(x_{ij}\) lies in the
corner with indices \(o(i),o(j)\). The matrix multiplication
rule follows by summing over the middle index and using
\(\sum_i e_i=1\). With one block this is the full algebra
\(M_n(B)\).

With two blocks with references \(p,q\), let \(M=e_pEe_q\).
Use the involution to identify \(e_qEe_p\) with \(M\).
For \(u,v\in M\), define \(\beta(u,v)\in B\) by
\(uv^*=\beta(u,v)e_p\). The involution fixes \(Be_p\),
so \(\beta\) is symmetric. Cyclicity of the trace shows that
\(u^*v=\beta(u,v)e_q\). Associativity in \(E\) gives
\(\beta(u,v)w=u\beta(v,w)\) for \(w\in M\).
Thus the rule
\[
       (b,u)(b',u')=(bb'+\beta(u,u'),\,bu'+b'u)
\]
makes \(B\oplus M\) a finite commutative \(B\)-algebra,
with an injective scalar map.

All the \(\beta(u,v)\) lie
in \(\mathfrak m_B\) because the mixed corners reduce to zero.
The coordinate changes identify \(E\) with the full space of
matrices having scalar diagonal blocks and entries in \(M\)
in the off-block rectangles, with involution transpose.
 In particular the group matrices
span the entire pattern over \(B\). Their integral
\(B_0\)-span is finite and continuous in the
\(\mathfrak m_{B_0}\)-adic sense proved for \(E_0\).
\label{cor:ordinary-patching-matrices-6}

\proofstep{The determinant law in the normalized basis.}
After extending scalars to the ring \(B\oplus M\), the determinant law on the group matrices is therefore
\[
 D(g)=\det\bigl(R(g)\bigr)
\]
in the indicated matrix realization \(R\).
One way to see the equality for any scalar algebra is to add
independent variables \(X_i\) to the diagonal and eliminate
successive pivots. Each pivot is a scalar in the enlarged
base once the previous pivots have been inverted. It is
monic in its new \(X_i\), and hence this localization is
injective. For \(i\ne j\), conjugation by \(1+(T-1)e_i\) over the
Laurent polynomial ring proves
\[
 D(1+a e_{ij})=D(1+Ta e_{ij}).
\]
Both sides are polynomials in \(T\); the polynomial ring
injects into its Laurent localization over any base.

At
\(T=0\) this gives \(D(1+a e_{ij})=1\).
On the diagonal, multiplying the idempotent factors gives
\[
  D(\operatorname{diag}(a_1,\dots,a_n))
    =\prod_i D(1+(a_i-1)e_i)=\prod_i a_i .
\]
Here the idempotent lines had rank one by the specialized
characteristic polynomial identity. Elimination now identifies
the two determinant laws before setting the \(X_i\) to zero.
The identity includes singular matrices and all coefficient
characteristics at issue.  The relation quotient imposes the ordered flag, fixed inertia and auxiliary identities. The transpose rule gives the polarization identity.
\label{cor:ordinary-patching-matrices-7}

\paragraph{Full matrix algebras and nonsemisimple residue algebras.}
The off-diagonal products determine which block pattern survives specialization. Write
\(I_M=(\beta(M,M))\). For a prime \(Q\) of \(B\),
extend to a field residue of \(B\oplus M\) over \(Q\).
Such a field exists by finiteness and the injective scalar
map. If \(M\) has nonzero image in this field, the span
contains every matrix unit by rescaling a nonzero entry
from the full mixed span. If \(M\) maps to zero, the span
is block diagonal. The mixed image is detected by its products: 
\[
 M\longmapsto0\text{ in }k'\quad\Longleftrightarrow\quad I_M\subseteq Q,
\]
 where \(k'\) is the chosen residue field. Indeed each mixed square is given by \(\beta\), and a square-zero field element vanishes.

Let \(k'/\kappa(Q)\) be the chosen residue-field extension. The group matrices have trace products in \(\kappa(Q)\) because the trace comes from \(B\). Applying Lemma~\ref{lem:matrix-span-by-traces} to the two-block calculation gives
\[
 I_M\not\subseteq Q
 \quad\Longleftrightarrow\quad
 \exists\,g_1,\ldots,g_{n^2}:\ 
 \det\!\left(\operatorname{tr}(R(g_a)R(g_b))\right)_{a,b=1}^{n^2}\ne0
 \quad\text{in }\kappa(Q).
\]
The matrices themselves are evaluated over \(k'\).
 In the one-block case the span is full.

For the other criterion suppose the residue algebra obtained
from the indicated Cayley--Hamilton kernel presents a
nonsemisimple algebra over a field with coefficient kernel
\(Q\). If \(b\in\operatorname{Ann}_B M\) lies outside \(Q\),
then all the mixed corners vanish after scalar extension
to this field. The Cayley--Hamilton quotient then becomes
the product of its two full diagonal matrix algebras.
The scalar extension statement follows by taking the span
of the kernel generators in the field group algebra, since
the group algebra surjects onto \(E\) with that defining
ideal. Any quotient of this product is semisimple, a
contradiction. This proves \(\operatorname{Ann}_B M\subset Q\). Both criteria supply the inputs for the
completed local comparison in the next subsection.
\label{cor:ordinary-patching-matrices-8}
\end{proof}

\subsection{From the Selmer ranks to the determinant locus}
\label{sec:patching-cotangent}
\hypertarget{source-ordinary-patching-cotangent}{}

The matrix realization of Section~\ref{sec:patching-matrices} lets us identify
infinitesimal deformations. We now use their dimensions to compare primes in the
completed localization at a comparison point. Write
\(B\) for that local ring, \(F=k((t))\) for its residue
field and \(r_F:B\to F\) for the residue map. The scalar
section is a complete coefficient DVR \(\mathcal O\)
mapping to \(B\) and inducing an isomorphism on residue
fields. The realization over the relative coefficient ring has the following properties. The coefficient algebra of
the matrices is finite over \(B\), with underlying module
\(B\oplus M\). The entries in the diagonal blocks are in
\(B\), those in the off-block rectangles in \(M\).
The representation matrices span all matrices with the specified pattern over \(B\).

In the product the off-block terms have coefficients in
\[
 M\cdot M\subseteq\mathfrak m_B\subset B .
\]
Write \(I_M\) for the ideal generated by these products.
In the one-block case \(M=0\).
The matrices satisfy the ordered local flag product
identity and the polarization and inertia relations.
The underlying global representation specializes to the
curve point. We use the finite matrix lattice over the
unlocalized determinant ring, continuous in its maximal
ideal topology, and the scalar lifts from the curve
ring \(k[[t]]\) in the coefficient section. These are
the matrix and coefficient inputs of the calculation.
\label{cor:ordinary-patching-cotangent-1}

\paragraph{The limit cohomology.}
The estimates available on the curve are finite-depth estimates. The first step is to
identify what they bound after passage to the limit.

\begin{lemma}[Passing the finite-depth bounds to the marked limit]\label{lem:marked-limit-cohomology}
Use the fixed curve and the marked system of
Proposition~\ref{prop:marked-profinite-system}.
Let \(Z_{\rm ord,b}\) be the bounded first-order cocycles satisfying the
ordinary local conditions on the limit and \(B_{\rm b}\) their subspace of
bounded boundaries. Then
\[
       \dim_F(Z_{\rm ord,b}/B_{\rm b})\le D_0=nd+nq .
\]
For \(g+1\) even tuples on the limit there is a nonzero
\(F\)-linear combination admitting the paired odd lift.
Any such lift has zero odd value on the marked tame generators,
and its restriction in the odd directions is read in the base
space with the chosen basis \((\kappa_i)_{i=1}^g\).
Here bounded means that a common power of \(t\) makes the
cocycles and the matrices imposing their local conditions integral.
\end{lemma}

\begin{proof}
Use the marked limit of the auxiliary groups from
Section~\ref{sec:patching-marking}. The
marking gives surjections from one finitely generated
profinite group to the groups at the auxiliary levels
and at the limit. In each finite quotient of the marking
group the level kernels have the same images as the
limit kernel on an ultrafilter extending the tails.
The marking also identifies the coefficient action,
complex conjugation, and the local restriction maps
on any fixed finite quotient on that ultrafilter.

\proofstep{Obtaining the finite-level cochains from the marked limit.}
At depth \(m\) pull back the cocycles to the positive
subgroup of the marking group. Let \(d\) be the number of topological generators and \(H\) the finite semidirect product at the chosen depth. Evaluation on those generators gives
\[
 \#\operatorname{Hom}_{\rm cont}(G,H)\le |H|^d,
\]
proving finiteness for the affine representations considered here. Each such
cocycle factors through a finite quotient of the
marking group. Take the simultaneous matching event
for these quotients and the coefficient action.
Descent to one level is equivalent there to descent
to the limit. The local cocycle values agree under
the marked local maps. The same matrix imposes the ordinary condition on either side, and unramifiedness is preserved as well.
The matrices generating boundaries also give exactly
the same pulled-back boundaries.

For this depth \(m\), let \(E_m\) be the matching
event for all the required finite quotients. If
\(N_{m,1}\) and \(N_{m,2}\) are thresholds for the two
finite-level estimates, choose a level in
\[
                 E_m\cap\{N\ge\max(N_{m,1},N_{m,2})\}.
\]
It is nonempty because both sets belong to the ultrafilter.
The finite quotient bound and the lift with factor \(t^C\)
transfer by the matched equations. This choice may depend on
\(m\), while the multiplier \(t^C\) is fixed. \label{cor:ordinary-patching-cotangent-2}

Choose \(e\) integral first-order cocycles with independent
classes over \(F\). The equation that an \(A\)-linear
combination is the boundary of a skew matrix is a system
of linear equations in the coefficients and the entries of
that matrix. Its rows span a finite \(A\)-module, so one
finite matrix supplies the same system at every \(t^m\).
Let \(B\) be this matrix, with domain the direct sum of
the coefficient module \(A^e\) and the module of integral skew matrices,
and let \(p\) project onto \(A^e\). Over \(F\) a relation among the chosen classes has zero coefficient vector.

In terms of the projection, 
\[
 \ker(B\otimes F)\subseteq\ker(p\otimes F).
\]
 Thus \(p\otimes F\) factors through the image of \(B\otimes F\); multiplying the factor map by \(t^{c_0}\)
for some fixed \(c_0\ge0\) makes it integral. Equivalently,
the nonzero Smith divisors of \(B\) give such a common
clearing exponent. Clearing denominators in the factor map gives an integral matrix \(L\) satisfying
\[
 LB=t^{c_0}p .
\]
For a vector satisfying the equations modulo \(t^m\),
\[
 Bv=t^mw\quad\Longrightarrow\quad
 t^{c_0}p(v)=L Bv=t^m Lw\equiv0\pmod{t^m}.
\]
Thus the kernel of the map from \((A/t^m)^e\) into the
finite-depth quotient is contained in
\((t^{m-c_0}A/t^m)^e\) for \(m\ge c_0\), giving an image
of size at least \(Q^{e(m-c_0)}\).
Proposition~\ref{prop:marked-first-order-bound} bounds that
quotient by \(Q^{mD_0+C}\) after finite-depth comparison.
Taking logarithms to base \(Q\) in these two estimates, then dividing by \(m\), gives
\[
 e(m-c_0)\le mD_0+C\quad\Longrightarrow\quad e\le D_0\qquad(m\to\infty).
\]

For each \(m\) apply the obstruction bound of
Proposition~\ref{prop:marked-odd-obstruction-bound} to the
\(g+1\) tuples reduced modulo \(t^m\). This is a bound for
lifting cocycles together with matrices imposing ordinarity. The failure
of a lifted tuple to satisfy the cocycle equation is an odd
two-cocycle, and its local corrections are the trivializations
in \textup{(11)}. Changing either lift adds the corresponding
coboundary in that complex. Therefore a unit-coefficient
combination multiplied by the fixed \(t^C\) has a lift.
Only the resulting one-cocycle and its local matrices are
transferred between the finite groups, through the finite
affine representations just described.

At each fixed depth there are finitely many coefficient vectors,
lifted cocycles and local matrices. A diagonal subsequence gives a
compatible reduction at every depth, hence an integral lift on
the limit. Write \((u_i)\) for the compatible coefficient vector. Unit-primitivity survives passage to the limit: 
\[
 \min_i v_t(u_i)=0,\qquad (t^Cu_i)_i\ne0\quad\text{in }F^{g+1}.
\]
 This gives the required nonzero combination. For a tame generator \(\tau_j\) the even data factor
through the base, where its image is trivial. The relation
\([\phi_j,\tau_j]=1\) therefore gives
\[
  \operatorname{Ad}(r(\phi_j))\,z(\tau_j)=z(\tau_j)
\]
for its odd cocycle value. The two Frobenius blocks have disjoint
spectra, so their off-diagonal centralizer is zero. The cocycle
rule and continuity now kill all of the closed normal subgroup
generated by the \(\tau_j\). Consequently the lift descends in
these directions to the ordinary odd space of the base, with
basis \((\kappa_i)\).
\label{cor:ordinary-patching-cotangent-3}

\proofstep{Passing to the fraction field with bounded denominators.}
Over \(F\) the three conclusions apply to bounded
continuous cocycles, meaning that some power of
\(t\) makes the cocycle integral and continuous on
every truncation. Clear one power for all entries
and for the finitely many local matrices, then
apply the integral conclusion. The cochain identities and local matrix equations are linear over the
coefficient ring and that ring injects into \(F\).

\label{cor:ordinary-patching-cotangent-4}

\end{proof}

Put
\[
       J_2=\mathfrak m_{\mathcal O}B+\mathfrak m_B^2+I_M,
       \qquad T=\bigl(\mathfrak m_B/J_2\bigr)^\vee .
\]
An element \(\alpha\in T\) extends uniquely to a derivation
\(B\to F\) at \(r_F\) killing \(\mathcal O\) and \(I_M\).
Indeed every element of \(B\) differs from a lift from
\(\mathcal O\) by an element of \(\mathfrak m_B\) and
products of two elements in \(\mathfrak m_B\) give zero.
The same equivalence identifies
\(\bigl(M/\mathfrak m_B M\bigr)^\vee\)
with additive maps \(\mu:M\to F\) that satisfy
\(\mu(bm)=r_F(b)\mu(m)\).
Apply \(\alpha\) to the diagonal block entries and
\(\mu\) to the off-block entries of the representation
matrices and multiply each result on the right by
the inverse of the base matrix. Denote the resulting
cocycles by \(z_\alpha\) and \(z_\mu\).
They are even and odd, respectively.

\begin{proposition}[Detection of the even and odd cotangent spaces]\label{prop:marked-cotangent-detection}
Use the normalized completed matrix model of
Proposition~\ref{prop:marked-finite-ch-algebra}. The even
derivation \(\alpha\in T\) and the odd functional
\(\mu\in(M/\mathfrak m_BM)^\vee\) give bounded cocycles satisfying
the indicated local and global conditions. Passing to cohomology detects each input separately: 
\[
 [z_\alpha]=0\ \Longrightarrow\ \alpha=0,\qquad [z_\mu]=0\ \Longrightarrow\ \mu=0.
\]
 In the two-block case choose
\(m_j\in M\) representing the coefficient of \([\kappa_j]\)
in the odd detection map. Then \(m_1,\ldots,m_g\) generate \(M\).
\end{proposition}

\begin{proof}
\proofstep{The paired test with its local relations.}
To see the multiplicative and local assertions for
these cocycles, use the paired test algebra over \(F\) with nilpotent basis
\(a_i,b_i,v\), products \(a_i b_j=\delta_{ij}v\) and all other
products of two basis elements zero. Indices range over whichever
finite family is being tested. For the original odd basis,
\(r_1=(1+\sum_j a_j\kappa_j)r\). Its adjoint on the square-zero
ideal with even \(b_j\) and odd \(v\) coordinates is the action
in the lifting assertion. Suppose that
\(\Phi:M\to F\) is \(\mathcal O\)-linear via \(r_F\)
and kills \(J_2M\). For \(m\in M\) define
\[
       d_m(b)=\Phi(bm)-r_F(b)\Phi(m).
\]
The map \(d_m\) kills \(\mathcal O\) and \(I_M\).
It is a derivation by writing
\(b=b_{\mathcal O}+b_1\) with \(b_1\in\mathfrak m_B\)
and using \(\Phi(\mathfrak m_B^2M)=0\).

Moreover \(d_{bm}=r_F(b)d_m\).
In a basis \(\alpha_i\) of \(T\) write
\(d_m=\sum_i\mu_i(m)\alpha_i\).
Each \(\mu_i\) is linear via \(r_F\).
Now the map on the coefficient algebra,
\[
 b+m\longmapsto
 r_F(b)+\sum_i a_i\mu_i(m)+\sum_i b_i\alpha_i(b)+v\Phi(m),
\]
is a ring map. In a mixed product \(bm\), the defining
derivation identity gives
\[
 \Phi(bm)=r_F(b)\Phi(m)+d_m(b)
         =r_F(b)\Phi(m)+\sum_i\alpha_i(b)\mu_i(m).
\]
Bilinearity and the multiplication table recover the required coefficient without mixed terms:
\[
 \left(\sum_i x_i a_i\right)\left(\sum_j y_j b_j\right)=\sum_{i,j}x_i y_j\delta_{ij}v=\left(\sum_i x_i y_i\right)v .
\] Products of
two odd elements have zero image: in the source they belong
to \(I_M\), killed by \(r_F\) and each \(\alpha_i\), and in
the target the odd-odd products are zero.
The same map works with any family \(\mu_i,\alpha_i\)
satisfying this formula for \(d_m\). By setting all
odd or all even data to zero, one obtains the respective
individual square-zero deformations defining
\(z_\alpha\) and \(z_\mu\).  \label{cor:ordinary-patching-cotangent-5}

The tests are bounded in the required topology.
Here is the role of the finite lattice in this argument.
Denote the nilpotent ideal of the test ring by \(J\) and choose \(r\) with \(J^r=0\). For any \(x_i\in J\) one has 
\[
 x_1^{e_1}\cdots x_d^{e_d}=0\qquad(e_1+\cdots+e_d\ge r).
\]
 The image of the unlocalized
determinant ring is contained in a finite
\(k[[t]]\)-algebra in that target. Indeed the scalar lifts
give the whole curve ring as constants. The kernel
of specialization to the curve ring in the unlocalized
determinant ring is finitely generated. The images
of those generators are nilpotent. If \(x_1,\ldots,x_d\) are the generator images, the scalar span of the following finite list already contains their image algebra:
\[
 \{x_1^{a_1}\cdots x_d^{a_d}:a_i\ge0,\ a_1+\cdots+a_d<r\}.
\]

A finite module of matrix
entries and differences follows by taking the images
of the generators of the original matrix lattice.
The images of the entries reduce to integral scalars.
Adjoining these entries using their nilpotent
differences from the scalars again gives a finite
algebra. If a determinant-ring element lies in the
\((j+r)\)-th power of its maximal ideal, expand its
products of maximal ideal elements into scalar parts
and nilpotent parts. The scalar parts lie in \((t)\)
and any term with at least \(r\) nilpotent factors
vanishes. The image is thus inside \(t^j\) times
the finite algebra. Continuity of the original
group matrices in their matrix lattice gives
continuity in these lattices for the test. Any
\(F\)-linear coordinate functional on them requires
only one fixed denominator.
\label{cor:ordinary-patching-cotangent-6}

\proofstep{The ordered local eigenspaces in the paired test.}
For a representation and ordered characters over the
test ring the ordinary relation is
\[
              \prod_{i=1}^n(\rho(h_i)-\psi_i(h_i))=0
                  \qquad(h_i\in G_v).
\]
The factors occur in order from \(1\) to \(n\).
By multilinearity the same identity holds on the group
algebra. Choose a finite linear combination \(z\) of group
elements for which the reduced character values at \(z\)
are pairwise different. Such a combination exists over
the infinite field \(F\): for each pair some group
element distinguishes it, and finitely many proper
linear hyperplanes cannot cover the space spanned by
these elements. The character differences lift to
units. Setting every factor at \(z\) gives the split
polynomial identity on \(\rho(z)\). The Lagrange
idempotents \(p_i\) of this polynomial are orthogonal
and sum to one.

For \(i\ge j\), put \(h\) in slot \(k\) with
\(j\le k\le i\), and \(z\) in all other slots.
Projecting the ordered product to \(p_i\) on the left
and \(p_j\) on the right gives
\[
 \left(\prod_{a<k}(\psi_i(z)-\psi_a(z))\right)
     p_i(\rho(h)-\psi_k(h))p_j
 \left(\prod_{a>k}(\psi_j(z)-\psi_a(z))\right)=0.
\]
The first product omits \(i\), the last omits \(j\);
both are units. If \(i>j\), the scalar term has
\(p_i p_j=0\), giving upper triangularity. If \(i=j=k\),
the surviving diagonal entry is \(\psi_i(h)p_i\).

\proofstep{Transport of the ordered relations through a nilpotent change.}
For two such systems agreeing modulo a square-zero
ideal, their idempotents \(p_i,p_i'\) agree modulo
that ideal. The matrix \(D=\sum_i p_i'p_i\) reduces to \(1\).
Orthogonality in each of the two idempotent systems gives
\[
       p'_jD=p'_jp_j=Dp_j,\qquad D^{-1}=1-(D-1).
\]
The inverse formula uses the square-zero ideal. Thus \(D\) carries
each old eigenspace to its new counterpart over the same ring. In the original ordinary frame the
original projectors are upper triangular, and a
basis made of their images of the respective standard
basis vectors differs from that frame by an upper
triangular unipotent change. Therefore
\(D^{-1}\rho'(h)D\) is upper triangular in the original
frame with the new prescribed diagonal. Both projectors
and \(D\) belong to the matrix subalgebra of the
allowed coefficient pattern. This follows from their
polynomial expressions and from inversion of scalar
units by a finite series in the nilpotent ideal.

Write \(D=1-Y\) and \(\rho'(h)=(1+\delta(h))r_1(h)\), with
\(Y,\delta(h)\) in the square-zero ideal. Multiplying the three
matrices gives
\[
  D^{-1}\rho'(h)D\,r_1(h)^{-1}
     =1+\delta(h)-\bigl(r_1(h)Y r_1(h)^{-1}-Y\bigr).
\]
Thus conjugation subtracts the coboundary of \(Y\), leaving a matrix in the required ordinary subspace.
On the prescribed local torsion subgroup the characters have fixed value,
so the resulting diagonal cocycle there vanishes.

A square-zero difference \(z\) therefore satisfies the differentiated multiplication and polarization laws 
\[
 z(gh)=z(g)+\operatorname{Ad}(r_1(g))z(h),\qquad z(\sigma h)=-z(h)^{\mathsf T},
\]
 in the polarized frame. The inertia identity gives the unramified
condition away from \(\ell\). This proves the local and
global claims for \(z_\alpha,z_\mu\) and for the
square-zero difference from \(r_1\) when the test
representation has the same \(a_i\) coordinates.
\label{cor:ordinary-patching-cotangent-7}

\proofstep{Detection modulo boundaries.}
If \(z_\alpha\) is a boundary, conjugation by a matrix
congruent to one kills the first-order representation
change. The determinant coefficients are then constant.
The flag characters are also constant. The ordered product
identity and the separating element of the group algebra
just used recover each character by its own projector.
The auxiliary Frobenius roots are unchanged because
their reduced differences are nonzero and the roots
are uniquely lifted for the same polynomial. On the corresponding spectral projectors the tame matrix is a conjugate of the constant identity matrix.

In particular for every allowed change of frame \(C\),
\[
 C\,\mathrm{id}\,C^{-1}-\mathrm{id}=0,
\]
so these tame variables vanish. The scalar parameter is fixed by the coefficient
section. The determinant coefficients, flag characters,
auxiliary roots and tame variables, and the scalar
parameter topologically generate the unlocalized ring.
Agreement on these elements gives agreement on the ring
by the finite lattice argument above. Localization and
completion preserve this equality in the test ring.
If \(J\) is the nilpotent test ideal and \(J^r=0\), each local test map \(f\) satisfies 
\[
 f(\widehat{\mathfrak m}^{\,r})\subseteq J^r=0.
\]
 Agreement on the generators thus extends to the completion and forces \(\alpha=0\).  \label{cor:ordinary-patching-cotangent-8}

Suppose \(z_\mu\) is the boundary of \(Y\).
Undoing the multiplication by the inverse base matrix, the
odd functional on each representation matrix is its scalar
specialization commuted with \(Y\). Both sides are \(B\)-linear
via \(r_F\), so the identity also holds on their \(B\)-span.
By the matrix-model hypothesis this span contains every
off-block matrix with a single entry \(m\in M\). Its scalar
specialization is zero, so the identity gives \(\mu(m)=0\).
This proves odd detection.

Write the odd detection map as
\[
 [z_\mu]=\sum_{j=1}^g \mu(m_j)[\kappa_j].
\]
Such \(m_j\) exist because the coefficients in the fixed basis
are linear functionals on \((M/\mathfrak m_BM)^\vee\), and
this finite-dimensional vector space has the natural double
dual. If the reductions of all \(m_j\) failed to span, there
would be a nonzero functional vanishing on them, contradicting
detection. Apply Nakayama's lemma to get the asserted generators.
\label{cor:ordinary-patching-cotangent-9}

\end{proof}

\paragraph{The spaces of liftable directions.}
To recover information about individual relations, we first separate the directions
that lift together. Let \(U\subseteq T^g\) consist of the tuples
\((\alpha_j)\) whose even cocycles \((z_{\alpha_j})\)
lift in the module with even \(b_j\) and odd \(v\)
coordinates, for the adjoint action of \(r_1\).
It is a linear subspace, since the possible lifts
and their local correction matrices satisfy linear equations.
The obstruction test on the limit, formulated in
Lemma~\ref{lem:marked-limit-cohomology}, shows that any
\(g+1\) vectors in \(T^g/U\) are linearly dependent, hence
\(\operatorname{codim}U\le g\).

\begin{proposition}[Dimension constraints on the liftable directions]\label{prop:marked-independent-lift-directions}
Let \(U\subset T^g\) consist of tuples whose even cocycles
lift to the module with even \(b_j\) and odd \(v\) coordinates
for the adjoint action of \(r_1\). Put
\(W=\{\alpha:e_j\alpha\in U\text{ for all }j\}\). Then
\[
 \operatorname{codim}U\le g,\qquad \dim_F W+g^2\le D_0 .
\]
For each \(\Phi\in(M/J_2M)^\vee\), the derivative tuple
\((d_{m_j})_{j=1}^g\) of the paired test belongs to \(U\).
\end{proposition}

\begin{proof}
\proofstep{Testing liftability by functionals on the obstruction module.}
For each functional \(\Phi\) on \(M/J_2M\), use
the paired test above. Its odd cocycle in coordinate
\(a_i\) has class
\(\sum_j\mu_i(m_j)[\kappa_j]\). Conjugation by one plus
a sum of the \(a_i\) times skew odd boundary generators
arranges equality of cocycles in these coordinates.
This leaves the even \(b_i\) coordinates unchanged.
Write \(C_{ij}=\mu_i(m_j)\).
Replace the \(a_i\)-coordinates by \(g\) coordinates
whose prescribed cocycles are \(\kappa_j\), and
replace the \(b_i\)-coordinates by their transpose
contractions with \(C\). Keep the scalar and \(v\)
coordinates. The resulting entries define a
representation even if \(C\) is rectangular or singular.
For multiplication, a pair of odd coordinate vectors
under this change satisfies
\[
           (C a)\cdot b=a\cdot(C^t b),
\]
so the product has exactly the old mixed coefficient.

The constant and first-order coordinates multiply by
the same formulas as before. The same calculation on
the ordered products, the polarization, and the inertia
identities preserves each of them. In relations with
zero right-hand side the new \(a_j\) parts vanish from
the corresponding relation for \(r_1\), while the other
parts vanish by the contraction.  Divide the new
representation by \(r_1\) on the right and subtract
one. The difference lies in the square-zero ideal
\((b_j,v)\), hence it is a cocycle for the adjoint action
of \(r_1\). The local flag argument above applies.
Its even coordinates are precisely the cocycles
of the \(g\) derivatives \(d_{m_j}\).
Consequently
\[
                  (d_{m_j})_{j=1}^g\in U
                    \qquad\text{for all }\Phi\in(M/J_2M)^\vee .
\]
Here this dual uses the \(F\)-vector space structure
from \(\mathcal O\), including the \(\mathcal O\)-linearity
required of \(\Phi\).  \label{cor:ordinary-patching-cotangent-10}

\proofstep{The subspace of independently liftable coordinates.}
Let \(e_j\alpha\) be the tuple supported in position \(j\). The condition on all positions is the intersection of the inverse images
\[
 W=\bigcap_j e_j^{-1}(U)=\{\alpha\in T:e_j\alpha\in U\ \text{for all }j\}.
\] For each \(\alpha\in W\), choose a lift
\(\Psi_j\) of \(e_jz_\alpha\). Put \(w_j=[v]\Psi_j\), the mixed coefficient of the lift. With \(a_i b_j=\delta_{ij}v\), the candidate cochain and its test images are 
\[
 Z=z_\alpha+\sum_{i=1}^g a_iw_i,\qquad b_jZ=b_jz_\alpha+vw_j=\Psi_j.
\]
 The scalar part and all products by
\(b_j\) determine such a first-order matrix. Thus the
cocycle, polarization, and inertia identities follow
from the corresponding identities for \(z_\alpha\)
and the \(\Psi_j\).

For the local condition, compare the even coordinates
of a matrix imposing ordinarity on \(\Psi_j\) with the tuple supported in
position \(j\) whose entry imposes the same condition on \(z_\alpha\). If \(D_{ij}\) is the difference in position \(i\) and \(\mathfrak b\) denotes the upper triangular space, local compatibility says 
\[
 \operatorname{Ad}(r(h))D_{ij}-D_{ij}\in\mathfrak b\qquad(h\in D_v).
\]

Each \(D_{ij}\) is therefore upper triangular over \(F\),
by the lower-entry induction used for the local
invariants in Section~\ref{subsec:odd-obstruction}.
In the ordinary first-order frame subtract their
lifts from the \(b_i\) coordinates of the matrix for \(\Psi_j\). These lifts are upper triangular.
Their boundaries stay upper triangular with diagonal
zero and do not change the required local condition.
Let \(Y\) impose the ordinary condition on \(z_\alpha\). The corrected matrix \(Y_j\) satisfies \([b_i]Y_j=\delta_{ij}Y\), and the same multiplication rule gives the combined matrix 
\[
 Y_* =Y+\sum_i a_i[v]Y_i,\qquad b_jY_*=Y_j.
\]
  \label{cor:ordinary-patching-cotangent-11}

For a basis of \(W\) these first-order classes, together
with the \(g^2\) classes \(a_j\kappa_i\), are independent.
Indeed a relation modulo boundaries in the scalar
coordinates first gives a boundary among the
\(z_\alpha\). The even detector forces the scalar
coefficients to vanish. The scalar part of the
boundary generator then commutes with the full
block span and is skew, hence is zero. In the remaining coordinates any relation has the form 
\[
 \sum_i c_{ij}[\kappa_i]=0\quad (1\le j\le g).
\]
 Independence of the odd classes gives \(c_{ij}=0\) for all \(i,j\).
We have constructed \(\dim W+g^2\) independent bounded
first-order classes on the same limit. Apply the dimension
bound of Lemma~\ref{lem:marked-limit-cohomology}.
\label{cor:ordinary-patching-cotangent-12}

\end{proof}

\paragraph{Independent coefficients of the odd relations.}
The liftable directions now have a role in the presentation: they distinguish the
coefficients of the odd relations one coordinate at a time.

\begin{proposition}[A presentation with independent odd-relation coefficients]\label{prop:marked-cotangent-presentation}
Use the paired cotangent detector of
Proposition~\ref{prop:marked-cotangent-detection} and the
per-coordinate liftable space \(W\) of
Proposition~\ref{prop:marked-independent-lift-directions}.
The completed marked ring \(B\oplus M\) is a quotient of
\[
 P=\mathcal O[[X_1,\ldots,X_{D_0},Y_1,\ldots,Y_g]] .
\]
There is a \(g\)-by-\(g\) matrix \(Z\) in the maximal ideal
of \(P\) for which the kernel contains the column \(ZY\).
The \(g^2+g\) entries of \(Z,Y\) are independent in the
relative cotangent space of \(P\) over its residue field.
\end{proposition}

\begin{proof}
Choose a spanning family of \(\operatorname{codim}U\) rows of
\(U^\perp\) and pad it to \(g\) rows. Denote its entries by
\(\phi_{ij}\in T^\vee\). For a tangent vector \(\alpha\), the entries of this matrix test all the supported tuples:
\[
 \phi_{ij}(\alpha)=0\ (\text{all }i,j)\quad\Longleftrightarrow\quad e_j\alpha\in U\ (\text{all }j).
\] Thus their common
annihilator is \(W\). A complement to their span has
dimension \(\dim W\le D_0-g^2\).
Form a spanning list of \(T^\vee\) by first listing all
\(g^2\) entries in designated slots and then a basis of this
complement, padding the end to length \(D_0\). When an entry
is expressed by projection to its designated slot, the
\(g^2\) coefficient vectors in \(F^{D_0}\) are independent.
The entries themselves may be dependent in \(T^\vee\).
Lift the whole spanning list to \(x_1,\ldots,x_{D_0}\in
\mathfrak m_B\) using \(T^\vee=\mathfrak m_B/J_2\).

\proofstep{Independent coefficient lifts in the ambient ring.}
Lift each coefficient to \(\mathcal O\) and write
\(z_{ij}\in\mathfrak m_B\) for the resulting linear
combination of the \(x_a\). For every such \(\Phi\), vanishing of the row on the tuple \((d_{m_j})\) gives
\[
 \Phi\!\left(\sum_j z_{ij}m_j\right)=0 .
\]
The dual separates \(M/J_2M\); thus the relation
holds modulo \(J_2M\).
Express the error as a combination of the
\(m_j\) with coefficients in \(J_2\), and
subtract it from the row. This gives an exact
relation \(\sum_j z'_{ij}m_j=0\).

The image of \(I_M\) in \(B\oplus M\) lies in the
square of its maximal ideal, because it comes from
products of two odd elements there. Consequently
the corrections have zero relative cotangent
class in \(B\oplus M\).
The complete local algebra \(B\oplus M\) is a
quotient of
\[
               P=\mathcal O[[X_1,\ldots,X_{D_0},Y_1,\ldots,Y_g]]
\]
by sending \(X_a,Y_j\) to \(x_a,m_j\); the
surjectivity follows from completeness and these
generators modulo the square and coefficient ideal.
Lift each corrected relation coefficient to
\(P\) with the same linear term modulo the
coefficient uniformizer. The entries of the
resulting matrix \(Z\), together with the \(Y_j\),
are independent in the relative cotangent space
of \(P\). The kernel contains the column \(ZY\).
 \label{cor:ordinary-patching-cotangent-13}

\end{proof}

Use Proposition~\ref{prop:marked-cotangent-presentation}
in the local comparison for
Proposition~\ref{prop:marked-chain-support}. We now relate
the determinant locus of the matrix \(Z\) to prime
containment in the completed patch ring.

\paragraph{Comparison of primes.}
We can now use dimensions to compare the primes through which support may pass. At a
two-block point the allowed target primes also depend on the odd module.

\begin{proposition}[Prime containment on the local comparison]
\label{prop:marked-stable-prime-inclusion}
For primes \(Z_0,Q_0\) of the one-block normalized comparison ring \(B\),
\[
      \dim(B/Z_0)\ge D_0+1\quad\Longrightarrow\quad Z_0\subseteq Q_0 .
\]
On the two-block comparison, a stable prime \(Z_0\) of that
quotient dimension lies in each of the admissible stable or
thick target primes \(Q_0\). Use the odd module \(M\) and the paired relation
matrix of Proposition~\ref{prop:marked-cotangent-presentation}.
\end{proposition}

\begin{proof}
We prove the local implication needed to propagate
the support. The one-block case has \(M=0\).
The same cotangent generators make \(B\) a quotient
of the regular local ring in \(D_0\) variables over
\(\mathcal O\). A prime of \(B\) with quotient dimension
at least \(1+D_0\) has zero preimage there. It is therefore
contained in every prime of \(B\).

\proofstep{The two-block ideal and its stable locus.}
In the two-block case write \(J=(ZY)\) in \(P\) and
\(I=J+(\det Z)\). The ideal \(I\) has a unique least
prime. To prove this, for a fixed index \(c\) adjoin
a column \(\lambda\) of variables with \(\lambda_c=1\)
and impose
\[
              Z\lambda=0,\qquad
              Y_i=\lambda_iY_c\quad(i\ne c).
\]
At any prime above \(\mathfrak m_P\), the entries of
\(Z\) and \(Y\) are still part of a regular parameter
system in the localized polynomial ring. Indeed one
can append lifts of regular parameters of the polynomial
fiber to the original parameters.

The chart equations
are cotangent independent there, with pivots in
\(Z_{ic}\) and \(Y_i\) for \(i\ne c\).
The quotient by those equations is regular local,
hence a domain. Let \(K_c\) be the contraction of
the chart ideal from the localization at
\(\mathfrak m_P P[\lambda]\). It contains \(I\):
the chart gives \(Y=Y_c\lambda\). Since \(Z\lambda=0\), the
adjugate identity gives
\[
 0=\operatorname{adj}(Z) Z\lambda=(\det Z)\lambda,\qquad
 0=(\det Z)\lambda_c=\det Z .
\]

Under any local map from \(P\) into a local domain
with a chart solution, \(K_c\) maps to zero.
Indeed extend the map by the given \(\lambda\) and
localize at the preimage of the target maximal ideal.
The chart ideal remains prime in this localization
and lies inside the extension of
\(\mathfrak m_P P[\lambda]\). A multiplier outside
\(\mathfrak m_P P[\lambda]\) in the localization criterion for
\(K_c\) thus lies outside the chart ideal and can be
cancelled. The resulting element belongs to the
chart ideal and maps to zero. In the generic
localized chart all \(\lambda_i\) are units.
Since \(\lambda_d\) is invertible here, normalization in chart \(d\) is given by
\[
 \lambda' = \lambda_d^{-1}\lambda,\qquad \lambda'_d=1.
\]
This takes a solution in chart \(c\) to one in chart \(d\), proving \(K_d\subseteq K_c\).
Reversing the indices proves equality.

\proofstep{Dimension after specializing the relation matrix.}
For a prime \(Q'\supseteq I\), choose a valuation ring
of the fraction field of \(P/Q'\) dominating that local
domain. If the image of \(Y\) is nonzero, choose a
coordinate of least valuation and divide by it.
The resulting \(\lambda\) lies in the valuation ring
and gives a chart solution. If \(Y=0\) in the residue field, the determinant condition on the image of \(Z\) gives a projective solution to the kernel equation:
\[
 \det Z=0\quad\Longrightarrow\quad \exists\,\lambda\ne0\text{ with }Z\lambda=0 .
\] Normalize it in the
same way. In either case the common contraction
\(K_c\) lies in \(Q'\). This proves the least-prime
assertion.

The dimension of \(P/J\) is at most \(1+D_0\).
In fact adjoining the \(g^2\) equations
\(Z_{ij}=\delta_{ij}Y_i\) gives the same radical as
setting all entries of \(Z,Y\) equal to zero,
since \((ZY)_i\) becomes \(Y_i^2\).
Cotangent independence makes all \(g^2+g\) entries of
\(Z,Y\) part of regular parameters, so killing them leaves
dimension \(\dim P-g^2-g\). Write \(J'\) for the ideal
obtained by adjoining the \(g^2\) equations to \(J\).
The radical comparison and the dimension inequality give
\[
       \dim(P/J)\le\dim(P/J')+g^2=\dim P-g .
\]
This is the asserted \(1+D_0\) bound. A prime over \(J\)
with quotient dimension at least \(1+D_0\) avoiding
some \(Y_j\) must therefore be the least prime over \(I\).
It contains \(I\) by the adjugate identity, and
strict containment above the least prime would
lose dimension.

\proofstep{The comparison at the stable and thick primes.}
Now let \(Z_0\subset B\) be stable with quotient
dimension at least \(1+D_0\). Stability in the
two-block case implies \(I_M\not\subseteq Z_0\).
Lift \(Z_0\) to the finite algebra \(B\oplus M\)
and take its preimage in \(P\). Its quotient
dimension is preserved by integrality. Some
\(Y_j\) avoids this prime, since containing all
the odd generators would force it to contain
their products. Therefore the preimage is the
least prime over \(I\).

For the target prime \(Q_0\) the two sufficient conditions in the comparison read
\[
 I_M\not\subseteq Q_0\quad\text{(stable)},\qquad \operatorname{Ann}_B M\subseteq Q_0\quad\text{(thick)} .
\] The determinant of the
exact relation matrix \((z'_{ij})\) annihilates
\(M\), by the adjugate identity. It lies in
\(Q_0\) in the case with nonsemisimple residue algebra. In the stable case its
product with \(I_M\) is zero, so primality gives
the same inclusion. A prime over \(Q_0\) in
\(B\oplus M\) thus pulls back to a prime over \(I\)
in \(P\). It contains the preimage for \(Z_0\).
Contracting gives \(Z_0\subseteq Q_0\).
For \(g=0\), the odd module is zero by its generation
statement, so the two-block stability hypothesis
cannot occur.
They give the prime inclusion once the patched
support provides the dimension lower bound at
the indicated curve points. \label{cor:ordinary-patching-cotangent-14}
\end{proof}

\subsection{Specialization with nonsemisimple residue algebra}
\label{subsec:closed-link}
\hypertarget{source-ordinary-closed-link}{}

One condition is still needed before the local prime comparison can be used along the
chain. Proposition~\ref{prop:marked-chain-support} requires that the specialized Cayley--Hamilton algebra remain nonsemisimple after extension of the residue field.
Let \(D\) be the ordinary determinant ring at the level, \(S\) its
ordered splitting algebra after the coefficient choices above and
\(B_0\to S\) the augmentation of the ring on the marked group.
Write \(G_0\) for the base CM group and \(G_*\twoheadrightarrow G_0\)
for the corresponding surjection from the marked CM group. Let \(P\)
be a prime of \(S\). Assume that over a field at the contracted
prime of \(D\), the specialization of the ordinary
Cayley--Hamilton algebra is nonsemisimple.
In the application its nonsemisimple quotient acts on the nonsplit
extension constructed in Section~\ref{sec:linked-components}.

\begin{proposition}[Nonsemisimplicity after specialization of the completed algebra]
\label{prop:marked-thick-realization}
Let \(P\subset S\) be a prime of the ordered splitting algebra
at which the ordinary Cayley--Hamilton algebra is nonsemisimple
over a field at the contraction to \(D\). For every prime \(Q\)
of a completed point ring \(B\) contracting to the preimage of
\(P\) under \(B_0\to S\), the specialization of the closed marked
relation quotient is nonsemisimple. In particular the mixed module satisfies 
\[
 \operatorname{Ann}_B M\subseteq Q.
\]

\end{proposition}

\begin{proof}
\proofstep{Nonsemisimplicity after scalar extension.}
First, for a finite-dimensional algebra over a field, a nonzero
radical is a nonzero nilpotent two-sided ideal. Extension of fields
preserves that ideal by faithful flatness. Thus nonsemisimplicity
persists over a common field. Two maps out of a ring with the same
prime kernel induce embeddings of the fraction field of its prime
quotient. The tensor product of the two target fields over that
fraction field is nonzero, and quotient by a prime followed by
passage to fractions gives a common field. If \(C=A/J\) presents the group algebra quotient before scalar extension to \(k'\), right exactness gives the specialized presentation
\[
 C\otimes k'\simeq (A\otimes k')/\operatorname{im}(J\otimes k').
\] These facts
give a nonsemisimple specialization over a field above \(P\). Finiteness here
follows from the spanning set for the integral Cayley--Hamilton
algebra in Section~\ref{sec:patching-matrices}.

\proofstep{Inclusion of the closed kernels before extending scalars.}
We need an inclusion of kernels \emph{before} the next field
extension. In \(S[G_0]\) form the ideal \(J'\) generated by the
image of the closed ordinary kernel from \(D[G_0]\). Put
\(C'=S[G_0]/J'\). The finite set of group elements spanning the
original closed quotient spans \(C'\) over \(S\). Continuity of the marked group algebra at each depth supplies an open normal subgroup with
\[
 U_b\triangleleft G_0,\qquad [u]-1\in\mathfrak m_S^b C'\quad(u\in U_b).
\]
In fact continuity on the finite original lattice gives this
assertion over \(D\), and \(D\to S\) is local. Also \(B_0\to S\)
is local, as can be seen on the common coefficient residue field.

Under augmentation, the determinant and its involution pull back from
the base. Applying the group and coefficient maps to a
Cayley--Hamilton relation thus gives the corresponding relation in
\(C'\). This includes polynomial-law relations with variable
coefficients: the base-extended law on the Cayley--Hamilton algebra
satisfies its characteristic polynomial after every scalar extension.
The local ordered products and the old inertia relations map to their
base relations, including their images under polarization. At an auxiliary place let \(\lambda_1,\dots,\lambda_n\)
be the ordered roots in \(S\) and \(\phi\) the image of
Frobenius. The image of the ordered relation is
\[
            \prod_{i=1}^{n}(\phi-\lambda_i)=0
\]
by Cayley--Hamilton for the actual characteristic
polynomial, keeping multiplicities. The tame generator
maps to \(1\) and the tame variables to zero, so the
image of each relation with a tame factor is also zero.
The involution gives the same for their polarized images.

\proofstep{The closed relation ideal and its specializations.}
Take now a member of the \emph{closure} of the marked relation ideal.
For each \(b\), the map to \(C'/\mathfrak m_S^bC'\)
kills the base image of \(U_b-1\) and the coefficient ideal
of depth \(b\). Therefore it factors through the group ring
on the corresponding finite quotient of \(G_*\) at that
depth. It kills each defining relation there by the preceding
calculation. For an element \(x\) of the closed ideal the induced map \(f\) therefore satisfies 
\[
 f(x)\bmod\mathfrak m_S^b C'=0\qquad(b\ge1).
\]
 The image lies in every power of the maximal ideal. Since \(C'\) is finite over the Noetherian local ring \(S\), Krull intersection gives
\[
 \bigcap_{b\ge1}\mathfrak m_S^b C'=0 .
\] This is the required inclusion of closed relation
kernels before extending the residue field.

Over the extension of \(\kappa(P)\) on which the base specialization is nonsemisimple, the surjection of groups induces
a surjection from the marked specialized group algebra modulo its
closed relations onto the base specialized algebra. Write \(E_{\rm mark}\twoheadrightarrow E_{\rm base}\) for the resulting quotient. The quotient of a semisimple finite-dimensional algebra is semisimple, so its contrapositive gives 
\[
 E_{\rm base}\text{ nonsemisimple}\quad\Longrightarrow\quad E_{\rm mark}\text{ nonsemisimple}.
\]
 

The common-field argument, this time over
\(B_0\) and with the fraction field at \(Q\), produces a field over which
the specialized algebra has a nonsemisimple quotient. The map from
\(B\) to this field has kernel \(Q\).
The completed group-algebra model of
Section~\ref{sec:patching-matrices} is the scalar extension of the
finite algebra with this closed presentation. It follows that
\(\operatorname{Ann}_B M\subseteq Q\) by the mixed-corner criterion
there. This proves the hypothesis on thick targets in the incidence
calculation, for links as well as the last target. The corresponding
statement for stable targets follows directly from the trace Gram
determinant after field extension. \label{cor:ordinary-closed-link-1}
\end{proof}


\subsection{Recovering the augmentation support}
\label{subsec:marked-augmentation}

We have compared primes above each curve point. To obtain a statement about ordinary
Hecke support, we must now follow them back to the common level by augmentation.

\begin{proposition}[Support transport across a marked chain]\label{prop:marked-chain-support}
Fix a finite chain of normalized curve points at the same
ordinary level, with the auxiliary system of
Proposition~\ref{prop:marked-profinite-system} and the local
comparison of Proposition~\ref{prop:marked-stable-prime-inclusion}. Suppose there is a stable prime of the
ordinary Hecke support below the initial curve kernel.
If successive curves share a thick prime as in the
completed local comparison, write \(Q_{\rm end}\) for the last stable or thick target prime
and \(J_{\rm ord}\) for the image of the ordinary Hecke ideal at the target. Then
\[
                      J_{\rm ord}\subseteq Q_{\rm end}.
\] The augmentation comparison is made at the
target, where the tame variables vanish.
\end{proposition}

\begin{proof}
\proofstep{The common augmented comparison.}
Use the ordinary finite-level modules of
Section~\ref{par:ordinary-patching} with the same
choice of auxiliary places for the finite chain.
Write \(J\subset D\) for the annihilator-support
ideal in the base determinant ring, and let
\(D^{\rm sp}\) be the finite free splitting algebra
adjoining the ordered Frobenius roots. A common
coefficient extension and the variable specializing
to \(t\) are included in the base weight algebra.
They are also included in the base module \(X\)
after extension to \(D^{\rm sp}\). For every prime
of that algebra its support is the inverse image
of \(V(J)\). Indeed, let \(D\to S\) be the base change, \(\mathfrak q\in\operatorname{Spec}S\) and \(\mathfrak p=\mathfrak q\cap D\). The fibre of the finite module is
\[
 (X\otimes_D S)\otimes_S k(\mathfrak q)
  \simeq (X\otimes_D k(\mathfrak p))
                    \otimes_{k(\mathfrak p)} k(\mathfrak q).
\]
A field extension preserves nonzero vector spaces. Applying Nakayama at each of the two primes identifies the support test with the ideal test:
\[
 \mathfrak q\in\operatorname{Supp}(X\otimes_D S)\quad\Longleftrightarrow\quad J\subseteq\mathfrak p .
\]

\proofstep{The fixed marked limit for the chosen finite chain.}
The common finitely generated limit group gives
a coefficient ring \(R_0\) acting on these modules
at each precision. Its augmentation maps to the
base splitting algebra. The limiting relation ideal
\(I\) includes the group, polarization, ordered
ordinary, old inertia and auxiliary root and inertia
relations. Write \(\mathbf y\) for the \(nq\) tame
parameters and \(f\) for the specializations of the
two-map product. At every curve point \(f\) is a unit
after localization. Its factors reduce there to
differences of distinct roots since the new norms
tend to one.

At each fixed tame degree and coefficient precision
the modules at sufficiently deep levels are free
over the truncated coefficient and tame algebra.
Let \(u\) and \(v\) be the limiting comparison maps between the augmentation of the patched module and the base module \(X\), in opposite directions. Their composites specialize to multiplication by the same \(f\), so on inverting \(f\),
\[
 (f^{-1}v)u=1,\qquad u(f^{-1}v)=1 .
\]
This identifies the augmentation after localization with \(X[1/f]\).
The tame graded pieces are the corresponding
free monomial pieces over that augmentation.
This is the same module construction as for
one arc.

In particular approximation of a power
of each relation by the finite-level annihilators
makes that relation nilpotent on the patched
module. Here both the tame degree and coefficient
precision are fixed before taking the auxiliary
level depth. One may quotient by the relations
for the support comparison. More explicitly,
if \(U\) is the finite patched module over the
completed localized ring and
\(I\subseteq\sqrt{\operatorname{Ann}U}\), the
finite module \(U/IU\) has the same support by
Nakayama. Equivalently the comparison may be
made at each support prime before imposing the
relations. At augmentation primes the identity
\[
  \mathfrak q\in\operatorname{Supp}U
  \quad\Longleftrightarrow\quad
  \mathfrak q/(\mathbf y)\in
       \operatorname{Supp}(U/(\mathbf y)U)
 \tag{16}\label{eq:local-patch-16}
\]
follows by local Nakayama at a prime containing
\((\mathbf y)\). Apply \eqref{eq:local-patch-16} to the original patched
module, whose augmentation is \(X[1/f]\); this
avoids requiring its quotient by \(I\) to have
the original graded pieces.

\proofstep{Support after completed localization.}
For completed localization at the \(i\)-th
curve kernel, the completed weight ring has
dimension \(1+nd\). Indeed its coefficient and
weight algebra before localization includes
the uniformizer, one parameter specializing
to \(t\), and \(nd\) weight parameters. The
evaluation quotient has dimension one and
the localized weight ring is regular. The
selected completed summand of the augmentation
module is finite free over this completed weight
ring. On completing the finite algebra acting faithfully on \(X\), let \(e_i\) be its mutually orthogonal local idempotents. They exhibit the summands by
\[
 1=\sum_i e_i,\qquad \widehat X=\bigoplus_i e_i\widehat X .
\]
The factor module is a direct summand of a finite
free module over a local ring, hence is free.
Quasi-regularity of \(\mathbf y\) means that the
graded pieces of the completed patched module are
the monomial pieces on this free augmentation.

Write \(W\) for the completed weight ring, \(Y=(y_1,\ldots,y_{nq})\), and \(U\) for the completed patched summand. Lift a \(W\)-basis of \(U/YU\), of size \(d_0\), to \(U\). The resulting map
\[
  W[[\mathbf y]]^{d_0}\longrightarrow U
\]
induces the isomorphism on each graded piece
\[
 \bigoplus_{|\alpha|=h}W^{d_0}\mathbf y^\alpha
      \xrightarrow{\ \sim\ }Y^hU/Y^{h+1}U\qquad(h\ge0)
\]
by quasi-regularity. Induction on \(h\) gives the isomorphism modulo \(Y^{h+1}\), and completeness gives it on the full modules. The support
algebra embeds in the endomorphisms of this
finite free module. It is finite torsion-free over
the regular power series ring, of dimension
\(1+nd+nq\).

Let \(S\) be the regular power series ring and \(C\) this finite support algebra. Fix a minimal prime \(P\) of \(C\). If \(0\ne s\in P\cap S\), its image in the zero-dimensional Noetherian local ring \(C_P\) is nilpotent. Clearing a denominator gives
\[
               u s^a=0\quad\text{for some }u\in C\setminus P,\ a\ge1 .
\]
Multiplication by \(s\) on the free module is injective and hence is injective on its endomorphisms and on \(C\). This would force \(u=0\). Therefore \(S\hookrightarrow C/P\) is finite and integral. Going-up lifts prime chains of \(S\), and incomparability makes the contractions of strict chains in \(C/P\) strict. It follows that
\[
 P\cap S=0,\qquad \dim C/P=\dim S=1+nd+nq.
\]

The same statement holds for the minimal support
primes after the relations are imposed, since
the underlying support is unchanged.

\proofstep{Movement across a link and augmentation at the target.}
For clarity, the support is moved across a link
as follows. Take the stable start prime
on the augmentation and lift it to the completion
at the starting point. It lies in the patched
support by \eqref{eq:local-patch-16}. Choose a minimal support prime
below it, which is stable because stability is
preserved under generization. Sections~\ref{sec:patching-cotangent}
and~\ref{subsec:closed-link} put this prime below a lift of the
outgoing thick prime. Let \(R\) be the common tame-adic patch ring. Imposing the tame parameters is precisely the quotient
\[
 \pi:R\twoheadrightarrow R/(\mathbf y),\qquad\ker\pi=(\mathbf y).
\] If \(\mathfrak t\) is the outgoing prime contracted to \(R\), then \((\mathbf y)\subseteq\mathfrak t\).

For the next curve kernel \(\mathfrak k'\) in \(R\), the link on the augmentation gives
\[
  \mathfrak p_{\rm st}\cap R\ \subseteq  \mathfrak t=\pi^{-1}(\mathfrak t/(\mathbf y))\ \subseteq\ \mathfrak k' ,
\]
where \(\mathfrak p_{\rm st}\) is the minimal support prime chosen in the starting completion. Lift this latter
contraction by faithful flatness to the next
completed localization. Its support membership
comes from base change of the common patched
module. The stable prime in this step need not
contain \((\mathbf y)\). After iterating, apply the
local comparison to the last target and use
\eqref{eq:local-patch-16} there, where the target does contain
\((\mathbf y)\). Contracting to the original
determinant ring gives the inclusion of \(J\)
at that target. This is the support step in
the finite-chain comparison. It uses the
finite-level relation and comparison maps at
the actual Hecke levels and the local
cotangent and dimension implication at each
point.
\hfill\mbox{\cite{milne-duality}}\label{cor:ordinary-marked-patching-1}
\end{proof}
